\chapter{Curvature}
\label{chap:pet-curvature}

Faithful transcription of Petersen, \emph{Riemannian Geometry} (GTM 171, 3rd ed.),
Chapter 3. Curvature is introduced as the failure of second covariant derivatives to
commute — the connection itself already plays the role of a gradient of the metric,
and curvature plays the role of its Laplacian/Hessian. From the curvature tensor the
chapter extracts the curvature operator, sectional curvature, Ricci curvature, and
scalar curvature, proves their basic symmetries (the two Bianchi identities and
Riemann's 1854 characterization of constant curvature), and records the coordinate
formula for curvature in terms of Christoffel symbols. The second part of the
chapter develops the \emph{fundamental equations of Riemannian geometry}: the
radial, tangential (Gauss), and normal (Codazzi--Mainardi) curvature equations for
a hypersurface cut out by a function $f$, specialized to distance functions,
Jacobi fields, and parallel fields, culminating in the qualitative theory of
conjugate and focal points. Much of the theory carries over unchanged to the
pseudo-Riemannian setting; only formulas involving contractions need modification.

\section{Curvature}

\subsection{The Curvature Tensor}

\begin{definition}[curvature tensor]
  \label{def:pet-ch3-curvature-tensor}
  \lean{PetersenLib.curvatureTensor}
  \leanok
  Let $(M,g)$ be a Riemannian manifold with Riemannian connection $\nabla$. The
  \emph{curvature tensor} is the $(1,3)$-tensor defined on vector fields $X,Y,Z$ by
  \[
    R(X,Y)Z=\nabla^2_{X,Y}Z-\nabla^2_{Y,X}Z
    =\nabla_X\nabla_YZ-\nabla_Y\nabla_XZ-\nabla_{[X,Y]}Z
    =[\nabla_X,\nabla_Y]Z-\nabla_{[X,Y]}Z.
  \]
  The first expression is also called the \emph{Ricci identity} and is written
  $R_{X,Y}Z=\nabla^2_{X,Y}Z-\nabla^2_{Y,X}Z$. The same formula defines the
  curvature of an arbitrary tensor $T$:
  \[
    R(X,Y)T=R_{X,Y}T=\nabla^2_{X,Y}T-\nabla^2_{Y,X}T.
  \]
  \source{petersen:page-0096}
\end{definition}

\begin{lemma}[tensoriality of the curvature tensor]
  \label{lem:pet-ch3-curvature-tensor-tensorial}
  \lean{PetersenLib.curvatureTensor_tensorial}
  \leanok
  \uses{def:pet-ch3-curvature-tensor}
  $R(X,Y)Z$ is $C^\infty(M)$-linear in $Z$ (as it already is in $X,Y$, being built
  from second covariant derivatives): for $f\in C^\infty(M)$,
  $R(X,Y)(fZ)=fR(X,Y)Z$.
  \source{petersen:page-0096}
\end{lemma}
\begin{proof}
  \uses{lem:pet-ch3-curvature-tensor-tensorial}
  \leanok
  Since both second covariant derivatives are already tensorial in $X$ and $Y$, it
  suffices to check tensoriality in $Z$. Using the product rule for the second
  covariant derivative,
  \begin{align*}
    R(X,Y)(fZ)&=\nabla^2_{X,Y}(fZ)-\nabla^2_{Y,X}(fZ)\\
    &=f\nabla^2_{X,Y}Z-f\nabla^2_{Y,X}Z\\
    &\quad+(\nabla^2_{X,Y}f)Z-(\nabla^2_{Y,X}f)Z\\
    &\quad+(\nabla_Yf)\nabla_XZ+(\nabla_Xf)\nabla_YZ\\
    &\quad-(\nabla_Xf)\nabla_YZ-(\nabla_Yf)\nabla_XZ\\
    &=f\big(\nabla^2_{X,Y}Z-\nabla^2_{Y,X}Z\big)=fR(X,Y)Z,
  \end{align*}
  where the two mixed terms cancel outright and the Hessian terms cancel because
  the Hessian of a function is symmetric, $\nabla^2_{X,Y}f=\nabla^2_{Y,X}f$.
\end{proof}

\begin{remark}[the Ricci identity as a derivation]
  \label{rem:pet-ch3-ricci-identity-derivation}
  \lean{PetersenLib.curvatureTensor_derivation}
  \leanok
  \uses{lem:pet-ch3-curvature-tensor-tensorial}
  Because the Hessian of a function is symmetric, $R_{X,Y}$ acts trivially on
  functions; this is the content of the Ricci identity
  $\nabla^2_{X,Y}-\nabla^2_{Y,X}=R_{X,Y}=R(X,Y)$, where on the right $R(X,Y)$ is
  regarded as a $(1,1)$-tensor acting on tensors as a derivation. In particular, if
  $T$ is a $(0,k)$-tensor,
  \[
    (R_{X,Y}T)(X_1,\dots,X_k)=-T(R(X,Y)X_1,\dots,X_k)-\dots-T(X_1,\dots,R(X,Y)X_k).
  \]
  \source{petersen:page-0097}
\end{remark}

\begin{definition}[$(0,4)$-curvature tensor]
  \label{def:pet-ch3-curvature-04-tensor}
  \lean{PetersenLib.curvatureTensorFour}
  \leanok
  \uses{def:pet-ch3-curvature-tensor}
  Using the metric $g$, the curvature tensor is changed to a $(0,4)$-tensor by
  \[
    R(X,Y,Z,W)=g(R(X,Y)Z,W).
  \]
  \source{petersen:page-0097}
\end{definition}

\begin{proposition}[Prop. 3.1.1 — symmetries of the curvature tensor]
  \label{prop:pet-ch3-curvature-symmetries}
  \lean{PetersenLib.curvatureTensor_symmetries}
  \leanok
  \uses{def:pet-ch3-curvature-04-tensor}
  The Riemannian curvature tensor $R(X,Y,Z,W)$ satisfies:
  \begin{enumerate}
    \item[(1)] skew-symmetry in the first two and in the last two entries:
    $R(X,Y,Z,W)=-R(Y,X,Z,W)=R(Y,X,W,Z)$;
    \item[(2)] symmetry between the first two and the last two entries:
    $R(X,Y,Z,W)=R(Z,W,X,Y)$;
    \item[(3)] Bianchi's first identity: $R(X,Y)Z+R(Z,X)Y+R(Y,Z)X=0$;
    \item[(4)] Bianchi's second identity:
    $(\nabla_ZR)_{X,Y}W+(\nabla_XR)_{Y,Z}W+(\nabla_YR)_{Z,X}W=0$.
  \end{enumerate}
  \source{petersen:page-0097,petersen:page-0098,petersen:page-0099,petersen:page-0100}
\end{proposition}
\begin{proof}
  \uses{prop:pet-ch3-curvature-symmetries}
  \leanok
  \textbf{(1), first half.} Skew-symmetry $R(X,Y,Z,W)=-R(Y,X,Z,W)$ holds because
  $R(X,Y)Z=-R(Y,X)Z$ directly from the definition of $R$.

  \textbf{(1), second half.} Recall $[X,Y]$ is defined implicitly by
  $D_XD_Yf-D_YD_Xf-D_{[X,Y]}f=0$ for functions $f$, i.e.\ $R(X,Y)f=0$. Using this
  on $f=\tfrac12g(Z,Z)$,
  \begin{align*}
    0&=R_{X,Y}\tfrac12g(Z,Z)\\
    &=\tfrac12D_XD_Yg(Z,Z)-\tfrac12D_YD_Xg(Z,Z)-\tfrac12D_{[X,Y]}g(Z,Z)\\
    &=D_Xg(\nabla_YZ,Z)-D_Yg(\nabla_XZ,Z)-g(\nabla_{[X,Y]}Z,Z)\\
    &=g(\nabla_X\nabla_YZ,Z)-g(\nabla_Y\nabla_XZ,Z)-g(\nabla_{[X,Y]}Z,Z)\\
    &\quad+g(\nabla_XZ,\nabla_YZ)-g(\nabla_XZ,\nabla_YZ)\\
    &=g\big(\nabla^2_{X,Y}Z,Z\big)-g\big(\nabla^2_{Y,X}Z,Z\big)=R(X,Y,Z,Z).
  \end{align*}
  Polarizing $R(X,Y,Z,Z)=0$ in $Z$ (replace $Z$ by $Z+W$):
  \[
    0=R(X,Y,Z+W,Z+W)=R(X,Y,Z,Z)+R(X,Y,W,W)+R(X,Y,Z,W)+R(X,Y,W,Z),
  \]
  so $R(X,Y,Z,W)=-R(X,Y,W,Z)$, i.e.\ $R(X,Y,Z,W)=R(Y,X,W,Z)$ using the first half.

  \textbf{(3).} Recall the Lie derivative of the connection,
  $(\mathcal L_X\nabla)_YZ=\mathcal L_X(\nabla_YZ)-\nabla_{\mathcal L_XY}Z-\nabla_Y\mathcal L_XZ$.
  Since $\mathcal L_XY=[X,Y]$ and $\mathcal L_XZ=[X,Z]$ for vector fields, and using
  $\nabla_XZ-\nabla_ZX=[X,Z]$ (torsion-freeness),
  \begin{align*}
    (\mathcal L_X\nabla)_YZ
    &=\nabla_X\nabla_YZ-\nabla_{\nabla_XY-\nabla_YX}Z-\nabla_Y(\nabla_XZ-\nabla_ZX)\\
    &=\nabla_X\nabla_YZ-\nabla_{\nabla_XY}Z-\nabla_Y\nabla_XZ
      +\nabla_{\nabla_YX}Z+\nabla_Y\nabla_ZX\\
    &=R_{X,Y}Z+\nabla^2_{Y,Z}X.
  \end{align*}
  The (algebraic) Jacobi identity for the Lie derivative of the Lie bracket,
  $(\mathcal L_X\mathcal L)_YZ=(\mathcal L_X\nabla)_YZ-(\mathcal L_X\nabla)_ZY=0$
  (this is the statement $\mathcal L_X[Y,Z]=[\mathcal L_XY,Z]+[Y,\mathcal L_XZ]$
  rewritten via torsion-freeness), gives
  \begin{align*}
    0&=(\mathcal L_X\nabla)_YZ-(\mathcal L_X\nabla)_ZY\\
    &=\big(R_{X,Y}Z+\nabla^2_{Y,Z}X\big)-\big(R_{X,Z}Y+\nabla^2_{Z,Y}X\big)\\
    &=R_{X,Y}Z-R_{X,Z}Y+\big(\nabla^2_{Y,Z}X-\nabla^2_{Z,Y}X\big)\\
    &=R_{X,Y}Z-R_{X,Z}Y+R_{Y,Z}X=R_{X,Y}Z+R_{Z,X}Y+R_{Y,Z}X,
  \end{align*}
  which is Bianchi's first identity.

  \textbf{(2).} This is a direct combinatorial consequence of (1) and (3). Applying
  (3) twice and (1) repeatedly,
  \begin{align*}
    R(X,Y,Z,W)&=-R(Z,X,Y,W)-R(Y,Z,X,W)\\
    &=R(Z,X,W,Y)+R(Y,Z,W,X)\\
    &=-R(W,Z,X,Y)-R(X,W,Z,Y)-R(W,Y,Z,X)-R(Z,W,Y,X)\\
    &=2R(Z,W,X,Y)+R(X,W,Y,Z)+R(W,Y,X,Z)\\
    &=2R(Z,W,X,Y)-R(Y,X,W,Z)=2R(Z,W,X,Y)-R(X,Y,Z,W),
  \end{align*}
  where the fourth line applies (3) to the triple $(X,W,Y)$ inside
  $R(X,W,Z,Y)+R(W,Y,Z,X)$ together with (1), and the last two lines use (1)
  again. Hence $2R(X,Y,Z,W)=2R(Z,W,X,Y)$.

  \textbf{(4).} The covariant derivative of the $(1,3)$-tensor $R$ is given by
  the Leibniz rule
  \[
    (\nabla_XR)_{Y,Z}W
    =\nabla_X\big(R(Y,Z)W\big)-R(\nabla_XY,Z)W-R(Y,\nabla_XZ)W-R(Y,Z)\nabla_XW.
  \]
  Expand every curvature term in the cyclic sum
  $(\nabla_XR)_{Y,Z}W+(\nabla_YR)_{Z,X}W+(\nabla_ZR)_{X,Y}W$ by the definition
  $R(A,B)C=\nabla_A\nabla_BC-\nabla_B\nabla_AC-\nabla_{[A,B]}C$. Four kinds of
  terms appear, and each kind cancels:
  \begin{itemize}
    \item[(i)] the twelve third-order terms $\nabla_U\nabla_V\nabla_SW$ cancel
    in pairs across the cyclic sum: for instance
    $+\nabla_X\nabla_Y(\nabla_ZW)$ arising from $\nabla_X\big(R(Y,Z)W\big)$
    cancels against $-\nabla_X\nabla_Y(\nabla_ZW)$ arising from
    $-R(X,Y)\nabla_ZW$ inside $(\nabla_ZR)_{X,Y}W$;
    \item[(ii)] by torsion-freeness $\nabla_XY-\nabla_YX=[X,Y]$ and linearity
    of $\nabla$ in its direction, the direction-slot corrections combine as
    $-\nabla_{\nabla_XY}(\nabla_ZW)+\nabla_{\nabla_YX}(\nabla_ZW)
    =-\nabla_{[X,Y]}(\nabla_ZW)$, which cancels the term
    $+\nabla_{[X,Y]}(\nabla_ZW)$ produced by $-R(X,Y)\nabla_ZW$; cyclically for
    the other two pairs;
    \item[(iii)] likewise the field-slot corrections combine as
    $\nabla_Z(\nabla_{\nabla_XY}W)-\nabla_Z(\nabla_{\nabla_YX}W)
    =\nabla_Z(\nabla_{[X,Y]}W)$, which cancels the term
    $-\nabla_Z(\nabla_{[X,Y]}W)$ coming from $\nabla_Z\big(R(X,Y)W\big)$;
    cyclically for the other two pairs;
    \item[(iv)] the remaining bracket-direction terms pair up, using
    torsion-freeness and antisymmetry of the bracket, as
    $\nabla_{[\nabla_XY,Z]}W+\nabla_{[Z,\nabla_YX]}W
    =\nabla_{[\nabla_XY-\nabla_YX,\,Z]}W=\nabla_{[[X,Y],Z]}W$ and cyclically,
    so their total is
    $\nabla_{[[X,Y],Z]+[[Y,Z],X]+[[Z,X],Y]}W=0$ by the Jacobi identity.
  \end{itemize}
  Hence the cyclic sum vanishes, which is Bianchi's second identity.
\end{proof}

\begin{example}[Example 3.1.2 — Euclidean space is flat]
  \label{ex:pet-ch3-euclidean-flat}
  \lean{PetersenLib.euclideanSpace_curvature_eq_zero}
  \leanok
  \uses{def:pet-ch3-curvature-tensor}
  $(\mathbb R^n,g_{\mathrm{eu}})$ has $R\equiv0$, since $\nabla_i\partial_j=0$ for the
  standard Cartesian coordinate vector fields.
  \source{petersen:page-0100}
\end{example}

\subsection{The Curvature Operator}

\begin{definition}[inner product on bivectors]
  \label{def:pet-ch3-bivector-inner-product}
  \lean{PetersenLib.bivectorInnerProduct}
  \leanok
  On the space $\Lambda^2T_pM$ of bivectors, with $e_1,\dots,e_n$ an orthonormal
  basis of $T_pM$, the inner product making $\{e_i\wedge e_j\}_{i<j}$ an
  orthonormal basis is denoted $g$ as well, and satisfies
  \[
    g(x\wedge y,v\wedge w)=g(x,v)g(y,w)-g(x,w)g(y,v)
    =\det\begin{pmatrix}g(x,v)&g(x,w)\\g(y,v)&g(y,w)\end{pmatrix}.
  \]
  \source{petersen:page-0100}
\end{definition}

\begin{definition}[bivectors as skew-symmetric maps]
  \label{def:pet-ch3-bivector-skew-map}
  \lean{PetersenLib.bivectorSkewMap}
  \leanok
  \uses{def:pet-ch3-bivector-inner-product}
  A bivector $x\wedge y$ is interpreted as the skew-symmetric transformation
  \[
    (x\wedge y)(v)=g(x,v)y-g(y,v)x,
  \]
  a counter-clockwise $90°$ rotation in $\operatorname{span}\{x,y\}$ when $x,y$ are
  orthonormal, satisfying
  $g(x\wedge y,v\wedge w)=g(x,v)g(y,w)-g(x,w)g(y,v)=g((x\wedge y)(v),w)$.
  \source{petersen:page-0100}
\end{definition}

\begin{lemma}[Jacobi identity for bivector maps]
  \label{lem:pet-ch3-bivector-jacobi-identity}
  \lean{PetersenLib.bivectorSkewMap_jacobi}
  \leanok
  \uses{def:pet-ch3-bivector-skew-map}
  $(x\wedge y)(z)+(y\wedge z)(x)+(z\wedge x)(y)=0$.
  \source{petersen:page-0100}
\end{lemma}
\begin{proof}
  \uses{lem:pet-ch3-bivector-jacobi-identity}
  \leanok
  Expand each term via the definition of the skew map:
  \[
    (x\wedge y)(z)=g(x,z)y-g(y,z)x,\quad
    (y\wedge z)(x)=g(y,x)z-g(z,x)y,\quad
    (z\wedge x)(y)=g(z,y)x-g(x,y)z.
  \]
  Summing and grouping by the vector factored out gives
  $\big(g(x,z)-g(z,x)\big)y+\big(g(y,x)-g(x,y)\big)z+\big(g(z,y)-g(y,z)\big)x=0$,
  since $g$ is symmetric.
\end{proof}

\begin{definition}[curvature operator]
  \label{def:pet-ch3-curvature-operator}
  \lean{PetersenLib.curvatureOperator}
  \leanok
  \uses{def:pet-ch3-curvature-04-tensor,def:pet-ch3-bivector-inner-product,def:pet-ch3-bivector-skew-map}
  By the symmetry properties of $R$, it defines a symmetric bilinear map
  $R:\Lambda^2TM\times\Lambda^2TM\to\mathbb R$,
  \[
    R\Big(\sum X_i\wedge Y_i,\sum V_j\wedge W_j\Big)=\sum R(X_i,Y_i,W_j,V_j)
  \]
  (note the reversal of $V_j,W_j$). Consequently there is a unique self-adjoint
  operator $\mathfrak R:\Lambda^2TM\to\Lambda^2TM$, the \emph{curvature operator},
  with
  \[
    g\Big(\mathfrak R\Big(\sum X_i\wedge Y_i\Big),\sum V_j\wedge W_j\Big)
    =\sum R(X_i,Y_i,W_j,V_j).
  \]
  \source{petersen:page-0101}
\end{definition}

\subsection{Sectional Curvature}

\begin{definition}[directional curvature operator]
  \label{def:pet-ch3-directional-curvature-operator}
  \lean{PetersenLib.directionalCurvatureOperator}
  \leanok
  \uses{def:pet-ch3-curvature-tensor,prop:pet-ch3-curvature-symmetries}
  For $v\in T_pM$, the \emph{directional} (or \emph{tidal force}) \emph{curvature
  operator} is $R_v(w)=R(w,v)v:T_pM\to T_pM$. It is self-adjoint and $v$ is always
  a zero eigenvector, by the symmetry properties of $R$.
  \source{petersen:page-0101}
\end{definition}

\begin{definition}[sectional curvature]
  \label{def:pet-ch3-sectional-curvature}
  \lean{PetersenLib.sectionalCurvature}
  \leanok
  \uses{def:pet-ch3-directional-curvature-operator,def:pet-ch3-bivector-inner-product}
  The \emph{sectional curvature} of $v,w\in T_pM$ is the normalized biquadratic form
  \[
    \sec(v,w)=\frac{g(R_v(w),w)}{g(v,v)g(w,w)-g(v,w)^2}
    =\frac{g(R(w,v)v,w)}{g(v\wedge w,v\wedge w)}.
  \]
  Since the denominator is the squared area of the parallelogram spanned by
  $v,w$, $\sec(v,w)$ depends only on the plane $\pi=\operatorname{span}\{v,w\}$.
  \source{petersen:page-0101}
\end{definition}

\begin{proposition}[Riemann, 1854]
  \label{prop:pet-ch3-riemann-constant-curvature-equivalence}
  \lean{PetersenLib.riemann_constantCurvature_equivalence}
  \leanok
  \uses{def:pet-ch3-sectional-curvature,def:pet-ch3-curvature-operator,prop:pet-ch3-curvature-symmetries,lem:pet-ch3-bivector-jacobi-identity}
  The following are equivalent, for fixed $p\in M$ and $k\in\mathbb R$:
  \begin{enumerate}
    \item[(1)] $\sec(\pi)=k$ for all $2$-planes $\pi\subset T_pM$;
    \item[(2)] $R(v_1,v_2)v_3=-k(v_1\wedge v_2)(v_3)$ for all $v_1,v_2,v_3\in T_pM$;
    \item[(3)] $R_v(w)=k\cdot(w-g(w,v)v)=k\cdot\operatorname{prj}_{v^\perp}(w)$ for all
    $w\in T_pM$ and unit $v$;
    \item[(4)] $\mathfrak R(\omega)=k\cdot\omega$ for all $\omega\in\Lambda^2T_pM$.
  \end{enumerate}
  \source{petersen:page-0101,petersen:page-0102,petersen:page-0103}
\end{proposition}
\begin{proof}
  \uses{prop:pet-ch3-riemann-constant-curvature-equivalence}
  \leanok
  (2) $\Rightarrow$ (3): setting $v_1=v$ unit, $v_2=w$, $v_3=v$ in (2) and using
  $(v\wedge w)(v)=g(v,v)w-g(w,v)v=w-g(w,v)v$ gives (3). (3) $\Rightarrow$ (1):
  for unit $v$, pairing (3) with $w$ gives the diagonal identity
  $g(R_v(w),w)=k\,(g(w,w)-g(w,v)^2)=k\,g(v\wedge w,v\wedge w)$. Both sides
  are homogeneous in $v$, so rescaling an arbitrary $v\ne0$ to a unit vector
  (and noting both sides vanish for $v=0$) extends the identity
  $R(x,y,x,y)=-k\,g(x\wedge y,x\wedge y)$ to all pairs $x,y$. For linearly
  independent $v,w$ the strict Cauchy--Schwarz inequality gives
  $g(v\wedge w,v\wedge w)>0$, and dividing yields $\sec(v,w)=k$, i.e.\ (1).

  (1) $\Rightarrow$ (2): introduce the multilinear maps
  \[
    R_k(v_1,v_2)v_3=-k(v_1\wedge v_2)(v_3),\qquad
    R_k(v_1,v_2,v_3,v_4)=-kg((v_1\wedge v_2)(v_3),v_4)=kg(v_1\wedge v_2,v_3\wedge v_4).
  \]
  These satisfy properties (1),(2),(3) of \cref{prop:pet-ch3-curvature-symmetries}
  (skew-symmetries directly from $g$ being symmetric and bilinear; the
  pair-symmetry from the explicit formula $kg(v_1\wedge v_2,v_3\wedge v_4)$; the
  first Bianchi identity from \cref{lem:pet-ch3-bivector-jacobi-identity}). Let
  $D(v_1,v_2,v_3,v_4)=R(v_1,v_2,v_3,v_4)-R_k(v_1,v_2,v_3,v_4)$; then $D$ also
  satisfies properties (1)--(3) of \cref{prop:pet-ch3-curvature-symmetries}, being
  a difference of two tensors that do. The hypothesis $\sec\equiv k$ gives
  $D(v,w,w,v)=0$ for all $v,w$ --- for linearly independent pairs by the
  hypothesis (the denominator $g(v\wedge w,v\wedge w)$ is positive by strict
  Cauchy--Schwarz), and for dependent pairs because both $R(v,w,w,v)$ and
  $g(v\wedge w,v\wedge w)$ vanish on degenerate pairs. Polarizing in $w=w_1+w_2$,
  \[
    0=D(v,w_1+w_2,w_1+w_2,v)=D(v,w_1,w_2,v)+D(v,w_2,w_1,v)=2D(v,w_1,w_2,v),
  \]
  using property (1) of $D$ to identify $D(v,w_1,w_1,v)=D(v,w_2,w_2,v)=0$ and to
  see $D(v,w_2,w_1,v)=D(v,w_1,w_2,v)$; hence
  $0=D(v,w_1,w_2,v)=-D(v,w_1,v,w_2)$ (skew-symmetry in the last two slots). Thus
  $D(X,Y,Z,W)=0$ whenever $W=X$; using properties (1),(2) of $D$ this forces $D$
  to be alternating in all four variables (each transposition of adjacent
  variables produces a sign, verified via the skew- and pair-symmetries), which
  together with Bianchi's first identity (property (3)) forces $D=0$: for an
  alternating $4$-tensor, $D(X,Y,Z)_{\text{cyc}}=0$ becomes $3D(X,Y,Z,W)=0$ for a
  fixed $W$. Hence $R=R_k$, which is (2).

  (2) $\Rightarrow$ (4): the operator $\omega\mapsto k\cdot\omega$ on
  $\Lambda^2T_pM$ is linear and self-adjoint, and by (2) it satisfies the
  defining property of \cref{def:pet-ch3-curvature-operator} on wedges:
  \begin{align*}
    g\big(k\cdot(x\wedge y),v\wedge w\big)
    =k\,g(x\wedge y,v\wedge w)
    =-k\,g\big((x\wedge y)(w),v\big)
    =R(x,y,w,v),
  \end{align*}
  using the antisymmetry $g(x\wedge y,w\wedge v)=-g(x\wedge y,v\wedge w)$.
  Since the wedges span $\Lambda^2T_pM$ (expand a skew map in an orthonormal
  basis), the operator with this defining property is unique, so
  $\mathfrak R(\omega)=k\cdot\omega$, giving (4).

  (4) $\Rightarrow$ (1): for any $v,w$, the defining property gives
  $g(\mathfrak R(v\wedge w),v\wedge w)=R(v,w,w,v)$, while (4) gives
  $g(k\cdot(v\wedge w),v\wedge w)=k\,g(v\wedge w,v\wedge w)$; hence the
  diagonal identity $R(v,w,v,w)=-k\,g(v\wedge w,v\wedge w)$ holds for all
  pairs, and dividing by $g(v\wedge w,v\wedge w)>0$ on linearly independent
  pairs gives (1) as before.
\end{proof}

\begin{definition}[constant curvature]
  \label{def:pet-ch3-constant-curvature}
  \lean{PetersenLib.HasConstantCurvature}
  \leanok
  \uses{prop:pet-ch3-riemann-constant-curvature-equivalence}
  $(M,g)$ has \emph{constant curvature} $k\in\mathbb R$ if one (equivalently, by
  \cref{prop:pet-ch3-riemann-constant-curvature-equivalence}, all) of the four
  equivalent conditions holds at every $p\in M$ with the same $k$.
  \source{petersen:page-0103}
\end{definition}

\subsection{Ricci Curvature}

\begin{definition}[Ricci curvature]
  \label{def:pet-ch3-ricci-curvature}
  \lean{PetersenLib.RicciCurvature}
  \leanok
  \uses{def:pet-ch3-curvature-tensor}
  For an orthonormal basis $e_1,\dots,e_n$ of $T_pM$, the \emph{Ricci curvature} is
  the trace
  \[
    \operatorname{Ric}(v,w)=\operatorname{tr}(x\mapsto R(x,v)w)
    =\sum_{i=1}^ng(R(e_i,v)w,e_i)=\sum_{i=1}^ng(R(v,e_i)e_i,w)
    =\sum_{i=1}^ng(R(e_i,w)v,e_i),
  \]
  a symmetric bilinear form, equivalently the symmetric $(1,1)$-tensor
  $\operatorname{Ric}(v)=\sum_{i=1}^nR(v,e_i)e_i$.
  \source{petersen:page-0103}
\end{definition}

\begin{definition}[Einstein manifold]
  \label{def:pet-ch3-einstein-manifold}
  \lean{PetersenLib.IsEinstein}
  \leanok
  \uses{def:pet-ch3-ricci-curvature}
  Write $\operatorname{Ric}\ge k$ if every eigenvalue of $\operatorname{Ric}(v)$ is $\ge
  k$, equivalently $\operatorname{Ric}(v,v)\ge kg(v,v)$ for all $v$. $(M,g)$ is an
  \emph{Einstein manifold with Einstein constant $k$} if $\operatorname{Ric}(v)=k\cdot v$
  for all $v$, equivalently $\operatorname{Ric}(v,w)=k\cdot g(v,w)$.
  \source{petersen:page-0103,petersen:page-0104}
\end{definition}

\begin{proposition}[Ricci curvature via sectional curvature]
  \label{prop:pet-ch3-ricci-sectional-relation}
  \lean{PetersenLib.ricciCurvature_eq_sum_sectionalCurvature}
  \leanok
  \uses{def:pet-ch3-ricci-curvature,def:pet-ch3-sectional-curvature}
  If $v\in T_pM$ is a unit vector completed to an orthonormal basis
  $\{v,e_2,\dots,e_n\}$, then
  $\operatorname{Ric}(v,v)=\sum_{i=2}^n\sec(v,e_i)$.
  \source{petersen:page-0104}
\end{proposition}
\begin{proof}
  \uses{prop:pet-ch3-ricci-sectional-relation}
  \leanok
  From \cref{def:pet-ch3-ricci-curvature} with the orthonormal basis
  $\{v,e_2,\dots,e_n\}$,
  \[
    \operatorname{Ric}(v,v)=g(R(v,v)v,v)+\sum_{i=2}^ng(R(e_i,v)v,e_i)
    =0+\sum_{i=2}^n\sec(v,e_i),
  \]
  using $R(v,v)v=0$ and $\sec(v,e_i)=g(R(e_i,v)v,e_i)$ for orthonormal $v,e_i$
  (denominator $1$ in \cref{def:pet-ch3-sectional-curvature}).
\end{proof}

\begin{corollary}[Ricci determines sec in dimensions $2$ and $3$]
  \label{cor:pet-ch3-ricci-determines-sec-low-dim}
  \lean{PetersenLib.einstein_iff_hasConstantCurvature_of_finrank_three}
  \leanok
  \uses{prop:pet-ch3-ricci-sectional-relation,def:pet-ch3-einstein-manifold,def:pet-ch3-constant-curvature}
  For $n=2$, and also for $n=3$, $\operatorname{Ric}$ determines $\sec$
  completely: for $n=3$ and orthonormal $\{e_1,e_2,e_3\}$,
  \[
    \begin{pmatrix}1&0&1\\1&1&0\\0&1&1\end{pmatrix}
    \begin{pmatrix}\sec(e_1,e_2)\\\sec(e_2,e_3)\\\sec(e_1,e_3)\end{pmatrix}
    =\begin{pmatrix}\operatorname{Ric}(e_1,e_1)\\\operatorname{Ric}(e_2,e_2)\\\operatorname{Ric}(e_3,e_3)\end{pmatrix},
  \]
  a system with determinant $2$; hence a $3$-manifold is Einstein if and only
  if it has constant sectional curvature.
  \source{petersen:page-0104}
\end{corollary}
\begin{proof}
  \uses{cor:pet-ch3-ricci-determines-sec-low-dim}
  \leanok
  From \cref{prop:pet-ch3-ricci-sectional-relation}, for $n=3$ and
  orthonormal $\{e_1,e_2,e_3\}$,
  $\sec(e_1,e_2)+\sec(e_1,e_3)=\operatorname{Ric}(e_1,e_1)$,
  $\sec(e_1,e_2)+\sec(e_2,e_3)=\operatorname{Ric}(e_2,e_2)$,
  $\sec(e_1,e_3)+\sec(e_2,e_3)=\operatorname{Ric}(e_3,e_3)$, i.e.\ the displayed
  linear system, whose matrix has determinant $2\ne0$; hence $\sec$ is
  recoverable from $\operatorname{Ric}$, so $\operatorname{Ric}=cg$ for a constant $c$
  forces all $\sec(e_i,e_j)$ (and hence, by polarization/continuity over all
  planes) to equal the same constant.
\end{proof}

\begin{remark}[constant curvature is Einstein]
  \label{rem:pet-ch3-constant-curvature-einstein}
  \lean{PetersenLib.constantCurvature_isEinstein}
  \leanok
  \uses{def:pet-ch3-constant-curvature,def:pet-ch3-einstein-manifold,prop:pet-ch3-ricci-sectional-relation}
  If $(M,g)$ has constant curvature $k$, it is Einstein with Einstein constant
  $(n-1)k$: for unit $v$, \cref{prop:pet-ch3-ricci-sectional-relation} gives
  $\operatorname{Ric}(v,v)=\sum_{i=2}^n\sec(v,e_i)=(n-1)k$.
  \source{petersen:page-0104}
\end{remark}

\begin{remark}[Einstein metrics that are not of constant curvature]
  \label{rem:pet-ch3-einstein-not-constant-curvature-examples}
  \lean{PetersenLib.einstein_examples}
  \leanok
  % Lean: the rigorous logical claim of the remark — the dimensional obstruction.
  % By prop:pet-ch3-ricci-sectional-relation an Einstein 3-manifold has constant
  % curvature, so non-constant-curvature Einstein metrics require dimension ≥ 4.
  % The three illustrative example families (S^n×S^n, Fubini–Study CP^n,
  % generalized Schwarzschild) are not formalised (they need product/CP^n/
  % warped-product curvature infrastructure beyond the chapter's main line).
  \uses{def:pet-ch3-einstein-manifold,prop:pet-ch3-ricci-sectional-relation}
  Three basic types of Einstein metrics without constant curvature: the product
  metric $S^n(1)\times S^n(1)$ with Einstein constant $n-1$; the Fubini--Study
  metric on $\mathbb{CP}^n$ with Einstein constant $2n+2$; and the generalized
  Schwarzschild metric $dr^2+\phi^2(r)d\theta^2+\rho^2(r)ds^2_{n-2}$ on
  $\mathbb R^2\times S^{n-2}$ ($n\ge4$), a doubly warped product with Einstein
  constant $0$. By \cref{prop:pet-ch3-ricci-sectional-relation} the search for
  such metrics must start in dimension $\ge4$.
  (Only this dimensional obstruction is formalised, as the theorem that a
  dimension-$3$ Einstein metric has constant curvature so that non-constant
  Einstein metrics require dimension $\ge4$; the three example metrics
  themselves are illustrative and are not formalised, as they require
  product / $\mathbb{CP}^n$ / warped-product curvature machinery beyond this
  chapter's main line.)
  \source{petersen:page-0104}
\end{remark}

\subsection{Scalar Curvature}

\begin{definition}[scalar curvature]
  \label{def:pet-ch3-scalar-curvature}
  \lean{PetersenLib.scalarCurvature}
  \leanok
  \uses{def:pet-ch3-ricci-curvature,def:pet-ch3-curvature-operator}
  The \emph{scalar curvature} is $\operatorname{scal}=\operatorname{tr}(\operatorname{Ric})=2\operatorname{tr}\mathfrak R$,
  a function $\operatorname{scal}:M\to\mathbb R$.
  \source{petersen:page-0104,petersen:page-0105}
\end{definition}

\begin{remark}[formulas for the scalar curvature]
  \label{rem:pet-ch3-scalar-curvature-formulas}
  \lean{PetersenLib.scalarCurvature_formulas}
  \leanok
  \uses{def:pet-ch3-scalar-curvature}
  In an orthonormal basis $e_1,\dots,e_n$,
  \[
    \operatorname{scal}=\sum_{j}g(\operatorname{Ric}(e_j),e_j)
    =\sum_{i,j}g(R(e_i,e_j)e_j,e_i)=\sum_{i,j}g(\mathfrak R(e_i\wedge e_j),e_i\wedge e_j)
    =2\sum_{i<j}g(\mathfrak R(e_i\wedge e_j),e_i\wedge e_j)=2\sum_{i<j}\sec(e_i,e_j).
  \]
  When $n=2$, $\operatorname{scal}(p)=2\sec(T_pM)$.
  \source{petersen:page-0105}
\end{remark}

\begin{notation}[divergence/adjoint of a tensor]
  \label{notation:pet-ch3-divergence-adjoint}
  \lean{PetersenLib.divergenceAdjoint}
  \leanok
  For a $(0,k)$-tensor $T$ and an orthonormal frame $E_1,\dots,E_n$ on $(M,g)$, the
  adjoint of the covariant derivative is the $(0,k-1)$-tensor
  \[
    (\nabla^*T)(X_1,\dots,X_{k-1})=-\sum_{i=1}^n(\nabla_{E_i}T)(E_i,X_1,\dots,X_{k-1}),
  \]
  independent of the choice of orthonormal frame.
  \source{petersen:page-0105,petersen:page-0106}
\end{notation}

\begin{proposition}[Prop. 3.1.5 — the contracted Bianchi identity]
  \label{prop:pet-ch3-contracted-bianchi}
  \lean{PetersenLib.contractedBianchiIdentity}
  \leanok
  \uses{def:pet-ch3-scalar-curvature,def:pet-ch3-ricci-curvature,notation:pet-ch3-divergence-adjoint,prop:pet-ch3-curvature-symmetries}
  On any Riemannian manifold, $d\operatorname{tr}(\operatorname{Ric})=d\operatorname{scal}=-2\nabla^*\operatorname{Ric}$.
  \source{petersen:page-0106}
\end{proposition}
\begin{proof}
  \uses{prop:pet-ch3-contracted-bianchi}
  \leanok
  Fix $p\in M$ and $W\in T_pM$, and let $E_1,\dots,E_n$ be an orthonormal frame
  with $(\nabla_{E_i}E_j)(p)=0$ for all $i,j$ (e.g.\ obtained by parallel
  transport of an orthonormal basis of $T_pM$ along radial geodesics, or via
  normal coordinates). At $p$, $\operatorname{scal}=\sum_{i,j}R(E_i,E_j,E_j,E_i)$, and
  differentiating in direction $W$ (the extra terms from differentiating the
  frame vanish at $p$ since $(\nabla_{E_i}E_j)(p)=0$),
  \[
    d\operatorname{scal}(W)\big|_p=\sum_{i,j}(\nabla_WR)(E_i,E_j,E_j,E_i).
  \]
  Bianchi's second identity, applied as a $(0,4)$-tensor with $Z=W$, $X=E_i$,
  $Y=E_j$, $V=E_j$, $U=E_i$, gives
  $(\nabla_WR)(E_i,E_j,E_j,E_i)=-(\nabla_{E_i}R)(E_j,W,E_j,E_i)-(\nabla_{E_j}R)(W,E_i,E_j,E_i)$.
  Summing over $i,j$ and relabelling $i\leftrightarrow j$ in the second sum, then
  using that $\nabla_{E_i}R$ inherits the algebraic symmetries of $R$ from
  \cref{prop:pet-ch3-curvature-symmetries} (in particular
  $(\nabla_{E_i}R)(W,E_j,E_i,E_j)=(\nabla_{E_i}R)(E_j,W,E_j,E_i)$, using
  skew-symmetry in each pair), the two sums coincide, so
  \[
    d\operatorname{scal}(W)\big|_p=-2\sum_{i,j}(\nabla_{E_i}R)(E_j,W,E_j,E_i).
  \]
  For fixed $i$, $\sum_j(\nabla_{E_i}R)(E_j,W,E_j,E_i)$ is, by
  \cref{def:pet-ch3-ricci-curvature} and its formula
  $\operatorname{Ric}(W,E_i)=\sum_jg(R(E_j,W)E_i,E_j)=\sum_jR(E_j,W,E_i,E_j)$,
  equal to $-\operatorname{Ric}(W,E_i)=-\operatorname{Ric}(E_i,W)$ (via
  skew-symmetry $R(E_j,W,E_j,E_i)=-R(E_j,W,E_i,E_j)$), and since contractions
  commute with covariant differentiation this identity differentiates to
  $\sum_j(\nabla_{E_i}R)(E_j,W,E_j,E_i)=-(\nabla_{E_i}\operatorname{Ric})(E_i,W)$. Hence
  \[
    d\operatorname{scal}(W)\big|_p=-2\sum_i\big(-(\nabla_{E_i}\operatorname{Ric})(E_i,W)\big)
    =2\sum_i(\nabla_{E_i}\operatorname{Ric})(E_i,W)=-2(\nabla^*\operatorname{Ric})(W)\big|_p,
  \]
  by \cref{notation:pet-ch3-divergence-adjoint}. As $p,W$ were arbitrary this is
  $d\operatorname{scal}=-2\nabla^*\operatorname{Ric}$.
\end{proof}

\begin{lemma}[Lemma 3.1.4 (Schur, 1886)]
  \label{lem:pet-ch3-schur-lemma}
  \lean{PetersenLib.schurLemma}
  \leanok
  \uses{def:pet-ch3-sectional-curvature,def:pet-ch3-einstein-manifold,prop:pet-ch3-contracted-bianchi,prop:pet-ch3-riemann-constant-curvature-equivalence}
  Let $(M,g)$ have dimension $n\ge3$ and satisfy, for a function $f:M\to\mathbb
  R$, either
  \begin{enumerate}
    \item[(1)] $\sec(\pi)=f(p)$ for all $2$-planes $\pi\subset T_pM$, $p\in M$, or
    \item[(2)] $\operatorname{Ric}(v)=(n-1)f(p)v$ for all $v\in T_pM$, $p\in M$.
  \end{enumerate}
  Then $f$ is constant; the metric has constant curvature, respectively is
  Einstein.
  \source{petersen:page-0105,petersen:page-0106}
\end{lemma}
\begin{proof}
  \uses{lem:pet-ch3-schur-lemma}
  \leanok
  If (1) holds, \cref{prop:pet-ch3-riemann-constant-curvature-equivalence} at
  each $p$ (with $k=f(p)$) gives $R_v(w)=f(p)(w-g(w,v)v)$ for unit $v$, and
  summing over an orthonormal completion as in
  \cref{prop:pet-ch3-ricci-sectional-relation} gives
  $\operatorname{Ric}(v)=(n-1)f(p)v$, which is (2). So it suffices to treat (2).

  Since $\operatorname{Ric}(v)=(n-1)f(p)v$, taking the trace,
  $\operatorname{scal}=n(n-1)f$, so $d\operatorname{scal}=n(n-1)df$. On the other hand,
  writing $\operatorname{Ric}=(n-1)f\cdot g$ as a $(0,2)$-tensor and using the product rule
  for $\nabla$ together with $\nabla g=0$,
  \[
    -\nabla^*\operatorname{Ric}(X)=\sum_i(\nabla_{E_i}\operatorname{Ric})(E_i,X)
    =(n-1)\sum_i\big[(\nabla_{E_i}f)g(E_i,X)+f(\nabla_{E_i}g)(E_i,X)\big]
    =(n-1)\sum_idf(E_i)g(E_i,X)=(n-1)df(X),
  \]
  using $\sum_ig(E_i,X)E_i=X$ in an orthonormal frame. So
  $\nabla^*\operatorname{Ric}=-(n-1)df$, and \cref{prop:pet-ch3-contracted-bianchi} gives
  \[
    n(n-1)df=d\operatorname{scal}=-2\nabla^*\operatorname{Ric}=2(n-1)df.
  \]
  Hence $(n-1)(n-2)df=0$; since $n\ge3$, $n-1\ne0$, so $(n-2)df=0$, and $n\ge3$
  forces $df=0$, i.e.\ $f$ is constant.
\end{proof}

\begin{corollary}[Corollary 3.1.6]
  \label{cor:pet-ch3-einstein-ricci-scalar-characterization}
  \lean{PetersenLib.einstein_iff_ricci_eq_scal_div_n_metric}
  \leanok
  \uses{def:pet-ch3-einstein-manifold,def:pet-ch3-scalar-curvature,lem:pet-ch3-schur-lemma}
  An $n(>2)$-dimensional Riemannian manifold $(M,g)$ is Einstein if and only if
  $\operatorname{Ric}=\dfrac{\operatorname{scal}}{n}g$.
  \source{petersen:page-0106}
\end{corollary}
\begin{proof}
  \uses{cor:pet-ch3-einstein-ricci-scalar-characterization}
  \leanok
  If $(M,g)$ is Einstein with constant $k$, then $\operatorname{Ric}=kg$, and taking the
  trace, $\operatorname{scal}=nk$, so $k=\operatorname{scal}/n$ and $\operatorname{Ric}=(\operatorname{scal}/n)g$.
  Conversely, if $\operatorname{Ric}=(\operatorname{scal}/n)g$, set $f=\operatorname{scal}/(n(n-1))$, so
  $\operatorname{Ric}(v)=(n-1)f(p)v$ for all $v$; condition (2) of
  \cref{lem:pet-ch3-schur-lemma} gives $f$ constant (as $n>2$), hence
  $\operatorname{scal}=n(n-1)f$ is a constant and $\operatorname{Ric}=(\operatorname{scal}/n)g$ with constant
  $\operatorname{scal}/n$, i.e.\ $(M,g)$ is Einstein.
\end{proof}

\subsection{Curvature in Local Coordinates}

\begin{proposition}[curvature tensor in coordinates]
  \label{prop:pet-ch3-curvature-coordinate-formula}
  \lean{PetersenLib.curvatureTensor_coordinates}
  \leanok
  \uses{def:pet-ch3-curvature-tensor}
  In local coordinates, writing $X=X^i\partial_i$, $Y=Y^j\partial_j$,
  $Z=Z^k\partial_k$ and $R(\partial_i,\partial_j)\partial_k=R^l_{ijk}\partial_l$, so
  $R(X,Y)Z=X^iY^jZ^kR^l_{ijk}\partial_l$, the coefficients are
  \[
    R^l_{ijk}=\partial_i\Gamma^l_{jk}-\partial_j\Gamma^l_{ik}
    +\Gamma^s_{jk}\Gamma^l_{is}-\Gamma^s_{ik}\Gamma^l_{js}.
  \]
  At a point $p$ where $\Gamma^k_{ij}|_p=0$ this simplifies to
  $R^l_{ijk}|_p=\partial_i\Gamma^l_{jk}|_p-\partial_j\Gamma^l_{ik}|_p$. Using the
  formula for the Christoffel symbols in terms of $g$, $R^l_{ijk}$ depends on the
  metric coefficients $g_{ij}$ and their first two derivatives.
  \source{petersen:page-0107}
\end{proposition}
\begin{proof}
  \uses{prop:pet-ch3-curvature-coordinate-formula}
  \leanok
  Using $\nabla_{\partial_i}\partial_k=\Gamma^s_{ik}\partial_s$,
  \begin{align*}
    R^l_{ijk}\partial_l&=R(\partial_i,\partial_j)\partial_k
    =\nabla_{\partial_i}\nabla_{\partial_j}\partial_k-\nabla_{\partial_j}\nabla_{\partial_i}\partial_k\\
    &=\nabla_{\partial_i}\big(\Gamma^s_{jk}\partial_s\big)-\nabla_{\partial_j}\big(\Gamma^s_{ik}\partial_s\big)\\
    &=\partial_i(\Gamma^l_{jk})\partial_l+\Gamma^s_{jk}\nabla_{\partial_i}\partial_s
      -\partial_j(\Gamma^l_{ik})\partial_l-\Gamma^s_{ik}\nabla_{\partial_j}\partial_s\\
    &=\big(\partial_i\Gamma^l_{jk}-\partial_j\Gamma^l_{ik}\big)\partial_l
      +\Gamma^s_{jk}\Gamma^l_{is}\partial_l-\Gamma^s_{ik}\Gamma^l_{js}\partial_l\\
    &=\big(\partial_i\Gamma^l_{jk}-\partial_j\Gamma^l_{ik}+\Gamma^s_{jk}\Gamma^l_{is}-\Gamma^s_{ik}\Gamma^l_{js}\big)\partial_l,
  \end{align*}
  which is the claimed formula. The simplification at a point where all
  Christoffel symbols vanish is immediate by dropping the quadratic terms.
\end{proof}

\begin{remark}[Remark 3.1.7 — index notation for $\operatorname{Ric}$ and $\operatorname{scal}$]
  \label{rem:pet-ch3-ricci-scalar-index-notation}
  \lean{PetersenLib.ricciScalar_indexNotation}
  \leanok
  \uses{def:pet-ch3-ricci-curvature,def:pet-ch3-scalar-curvature,prop:pet-ch3-curvature-coordinate-formula}
  In coordinates one writes $\operatorname{Ric}_{ij}=R_{ij}=R^k_{ik}=g^{kl}R_{kijl}$ and
  $\operatorname{scal}=R=g^{ij}R_{ij}$, both obtained by contracting indices of the curvature
  tensor; to avoid confusing the full curvature tensor with the scalar curvature
  it is also denoted $\mathfrak{Rm}$.
  \source{petersen:page-0108}
\end{remark}

\begin{remark}[Remark 3.1.8 — lowering the last index of $R^l_{ijk}$]
  \label{rem:pet-ch3-curvature-index-conventions}
  \lean{PetersenLib.curvatureTensor_indexLowering}
  \leanok
  \uses{def:pet-ch3-curvature-04-tensor,prop:pet-ch3-curvature-coordinate-formula}
  Consistently with $R(X,Y,Z,W)=g(R(X,Y)Z,W)$, the fully lower-index curvature
  components are $R_{ijkl}=g_{sl}R^s_{ijk}=R(\partial_i,\partial_j,\partial_k,\partial_l)$.
  \source{petersen:page-0108}
\end{remark}

\section{The Equations of Riemannian Geometry}

\subsection{Curvature Equations}

\begin{definition}[the tensors $S$ and $\operatorname{Hess}^2f$]
  \label{def:pet-ch3-hessian-square-notation}
  \lean{PetersenLib.hessianOperator}
  \leanok
  For a smooth function $f$ (possibly only on an open $O\subset M$), write
  $S(X)=\nabla_X\nabla f$ for the $(1,1)$-tensor corresponding to
  $\operatorname{Hess}f$, and $\operatorname{Hess}^2f$ for the $(0,2)$-tensor
  corresponding to $S^2=S\circ S$.
  \source{petersen:page-0108}
\end{definition}

\begin{definition}[second fundamental form]
  \label{def:pet-ch3-second-fundamental-form}
  \lean{PetersenLib.secondFundamentalForm}
  \leanok
  For a hypersurface $H^{n-1}\subset M^n$ with fixed unit normal field
  $N:H\to T^\perp H=\{v\in T_pM\mid p\in H,\,v\perp T_pH\}$, the \emph{second
  fundamental form} is the $(0,2)$-tensor $\Pi(X,Y)=g(\nabla_XN,Y)$ on $H$.
  \source{petersen:page-0108,petersen:page-0109}
\end{definition}
\begin{proof}
  \uses{def:pet-ch3-second-fundamental-form}
  \leanok
  For $X,Y\in TH$, $[X,Y]\in TH$ as well, so all three are perpendicular to $N$.
  Since $\nabla$ is torsion-free, $\nabla_XY-\nabla_YX=[X,Y]\in TH$, and
  \[
    g(\nabla_XN,Y)=D_Xg(N,Y)-g(N,\nabla_XY)=-g(N,\nabla_XY)=-g(N,\nabla_YX)=g(\nabla_YN,X),
  \]
  so $\Pi$ is symmetric. Moreover $g(\nabla_XN,N)=\tfrac12D_X|N|^2=0$, so
  $\nabla_XN$ is automatically tangent to $H$ for $X\in TH$, and $\Pi(X,Y)$ is
  well defined by this formula for $X,Y\in TH$.
\end{proof}

\begin{proposition}[Proposition 3.2.1]
  \label{prop:pet-ch3-normal-hessian-relation}
  \lean{PetersenLib.normal_hessian_relation}
  \leanok
  \uses{def:pet-ch3-second-fundamental-form,def:pet-ch3-hessian-square-notation}
  Let $H\subset f^{-1}(a)$ be open and consist entirely of regular points of $f$.
  Then:
  \begin{enumerate}
    \item[(1)] $N=\nabla f/|\nabla f|$ is a unit normal to $H$;
    \item[(2)] $\Pi(X,Y)=\dfrac1{|\nabla f|}\operatorname{Hess}f(X,Y)$ for
    $X,Y\in TH$;
    \item[(3)] $\operatorname{Hess}f(\nabla f,X)=\tfrac12D_X|\nabla f|^2$ for
    $X\in TM$.
  \end{enumerate}
  \source{petersen:page-0108,petersen:page-0109}
\end{proposition}
\begin{proof}
  \uses{prop:pet-ch3-normal-hessian-relation}
  \leanok
  (1) $N=\nabla f/|\nabla f|$ has unit length by construction, and is
  perpendicular to $H$ since $D_Xf=0$ for $X$ tangent to $H$ (as $H\subset
  f^{-1}(a)$).

  (2) For $X,Y\in TH$,
  \begin{align*}
    \Pi(X,Y)&=g\Big(\nabla_X\frac{\nabla f}{|\nabla f|},Y\Big)
    =g\Big(\frac1{|\nabla f|}\nabla_X\nabla f,Y\Big)+g\Big(D_X\Big(\frac1{|\nabla
    f|}\Big)\nabla f,Y\Big)\\
    &=\frac1{|\nabla f|}\operatorname{Hess}f(X,Y),
  \end{align*}
  since $\nabla f\perp TH$ so $g(\nabla f,Y)=0$.

  (3) By symmetry of $\operatorname{Hess}f$,
  \[
    \operatorname{Hess}f(\nabla f,X)=g(\nabla_{\nabla f}\nabla f,X)
    =g(\nabla_X\nabla f,\nabla f)=\tfrac12D_X|\nabla f|^2.
  \]
\end{proof}

\begin{theorem}[Theorem 3.2.2 — the radial curvature equation]
  \label{thm:pet-ch3-radial-curvature-equation}
  \lean{PetersenLib.radialCurvatureEquation}
  \leanok
  \uses{def:pet-ch3-hessian-square-notation,def:pet-ch3-curvature-tensor}
  With $S,\operatorname{Hess}^2f$ as in \cref{def:pet-ch3-hessian-square-notation}:
  \begin{align}
    (\nabla_{\nabla f}S)(X)+S^2(X)-\nabla_X(S(\nabla f))&=-R(X,\nabla f)\nabla f,\\
    \nabla_{\nabla f}\operatorname{Hess}f+\operatorname{Hess}^2f-\operatorname{Hess}\Big(\tfrac12|\nabla
    f|^2\Big)&=-R(\cdot,\nabla f,\nabla f,\cdot),\\
    \mathcal L_{\nabla f}\operatorname{Hess}f-\operatorname{Hess}^2f-\operatorname{Hess}\Big(\tfrac12|\nabla
    f|^2\Big)&=-R(\cdot,\nabla f,\nabla f,\cdot).
  \end{align}
  \source{petersen:page-0110}
\end{theorem}
\begin{proof}
  \uses{thm:pet-ch3-radial-curvature-equation}
  \leanok
  For the first formula, using $\nabla_X\nabla f=S(X)$ and the definition of $R$,
  \begin{align*}
    -R(X,\nabla f)\nabla f&=-\nabla^2_{X,\nabla f}\nabla f+\nabla^2_{\nabla f,X}\nabla f\\
    &=-(\nabla_XS)(\nabla f)+(\nabla_{\nabla f}S)(X)\\
    &=-\nabla_X(S(\nabla f))+\nabla_{\nabla_X\nabla f}\nabla f+(\nabla_{\nabla f}S)(X)\\
    &=-\nabla_X(S(\nabla f))+S^2(X)+(\nabla_{\nabla f}S)(X),
  \end{align*}
  using the product rule $(\nabla_XS)(\nabla f)=\nabla_X(S(\nabla f))-S(\nabla_X\nabla
  f)=\nabla_X(S(\nabla f))-S(S(X))$. This is the first formula.

  For the second, since $\nabla(\tfrac12|\nabla f|^2)=\nabla_{\nabla f}\nabla f$ (the
  gradient of $\tfrac12|\nabla f|^2$), and covariant differentiation commutes
  with the type change from $(1,1)$- to $(0,2)$-tensors, so that
  $(\nabla_N\operatorname{Hess}f)(X,Y)=g((\nabla_NS)(X),Y)$, applying $g(-,Y)$ to
  the first formula (with $N=\nabla f$) and using
  $g(\nabla_X(S(\nabla f)),Y)=D_X(g(S(\nabla f),Y))-g(S(\nabla f),\nabla_XY)$ reduces
  precisely to the $g$-paired version of the first identity, namely
  $\nabla_{\nabla f}\operatorname{Hess}f+\operatorname{Hess}^2f-\operatorname{Hess}(\tfrac12|\nabla
  f|^2)=-R(\cdot,\nabla f,\nabla f,\cdot)$.

  For the third, compute the Lie derivative of $\operatorname{Hess}f$ along
  $\nabla f$ directly:
  \begin{align*}
    (\mathcal L_{\nabla f}\operatorname{Hess}f)(X,Y)
    &=(\nabla_{\nabla f}\operatorname{Hess}f)(X,Y)+\operatorname{Hess}f(\nabla_X\nabla
    f,Y)+\operatorname{Hess}f(X,\nabla_Y\nabla f)\\
    &=(\nabla_{\nabla f}\operatorname{Hess}f)(X,Y)+2\operatorname{Hess}^2f(X,Y),
  \end{align*}
  using $\operatorname{Hess}f(\nabla_X\nabla f,Y)=\operatorname{Hess}f(S(X),Y)=\operatorname{Hess}^2f(X,Y)$
  and symmetry of $\operatorname{Hess}^2f$. Substituting the second formula gives the
  third.
\end{proof}

\begin{remark}[Remark 3.2.3 — tangential/normal decomposition]
  \label{rem:pet-ch3-tangential-normal-decomposition}
  \lean{PetersenLib.tangentialNormalDecomposition}
  \leanok
  \uses{thm:pet-ch3-radial-curvature-equation}
  The third formula of \cref{thm:pet-ch3-radial-curvature-equation} computes a
  suitable curvature using only gradients of functions and Lie derivatives — no
  covariant derivatives are needed. The next two fundamental equations, known as
  the \emph{Gauss equation} and the \emph{Peterson--Codazzi--Mainardi equation},
  use the notation
  $X=X^\top+X^\perp=X-g(X,N)N+g(X,N)N$
  decomposing a vector into parts tangential and normal to a hypersurface $H$
  with normal $N$.
  \source{petersen:page-0110}
\end{remark}

\begin{theorem}[Theorem 3.2.4 — the tangential (Gauss) curvature equation]
  \label{thm:pet-ch3-tangential-curvature-equation}
  \lean{PetersenLib.tangentialCurvatureEquation}
  \leanok
  \uses{def:pet-ch3-second-fundamental-form,rem:pet-ch3-tangential-normal-decomposition}
  For $X,Y,Z,W$ tangent to a hypersurface $H$ with induced metric $g_H$ and
  induced curvature $R^H$,
  \[
    g(R(X,Y)Z,W)=g_H\big(R^H(X,Y)Z,W\big)-\Pi(X,W)\Pi(Y,Z)+\Pi(X,Z)\Pi(Y,W).
  \]
  \source{petersen:page-0111}
\end{theorem}
\begin{proof}
  \uses{thm:pet-ch3-tangential-curvature-equation}
  For $X,Y$ tangent to $H$, uniqueness of the Riemannian connection on $H$ gives
  $\nabla^H_XY=(\nabla_XY)^\top=\nabla_XY-g(\nabla_XY,N)N=\nabla_XY+\Pi(X,Y)N$
  (one checks $(\nabla_XY)^\top$ satisfies the defining properties of the
  Riemannian connection on $H$, e.g.\ via the Koszul formula). Using
  $\nabla_XY=\nabla^H_XY-\Pi(X,Y)N$, expand $R(X,Y)Z$ for $X,Y,Z$ tangent to $H$:
  \begin{align*}
    R(X,Y)Z&=\nabla_X\nabla_YZ-\nabla_Y\nabla_XZ-\nabla_{[X,Y]}Z\\
    &=\nabla_X\big(\nabla^H_YZ-\Pi(Y,Z)N\big)-\nabla_Y\big(\nabla^H_XZ-\Pi(X,Z)N\big)
      -\nabla^H_{[X,Y]}Z+\Pi([X,Y],Z)N\\
    &=\nabla_X\nabla^H_YZ-\nabla_Y\nabla^H_XZ-\nabla^H_{[X,Y]}Z\\
    &\quad-\nabla_X\big(\Pi(Y,Z)N\big)+\nabla_Y\big(\Pi(X,Z)N\big)+\Pi([X,Y],Z)N\\
    &=R^H(X,Y)Z-\Pi(X,\nabla^H_YZ)N+\Pi(Y,\nabla^H_XZ)N\\
    &\quad-\big(D_X\Pi(Y,Z)\big)N-\Pi(Y,Z)\nabla_XN+\big(D_Y\Pi(X,Z)\big)N+\Pi(X,Z)\nabla_YN\\
    &\quad+\Pi([X,Y],Z)N+\Pi(\nabla_XY,Z)N-\Pi(\nabla_YX,Z)N,
  \end{align*}
  using $[X,Y]=\nabla_XY-\nabla_YX$ and expanding
  $\nabla_X(\Pi(Y,Z)N)=(D_X\Pi(Y,Z))N+\Pi(Y,Z)\nabla_XN$. The terms
  $-\Pi(X,\nabla^H_YZ)N+\Pi(\nabla_XY,Z)N$ and
  $\Pi(Y,\nabla^H_XZ)N-\Pi(\nabla_YX,Z)N$ and $\Pi([X,Y],Z)N$ combine, using
  bilinearity of $\Pi$ and $\nabla_XY=\nabla^H_XY+\Pi(X,Y)N$, so that all
  quadratic-in-$\Pi$ tangential-index terms cancel, leaving
  \[
    R(X,Y)Z=R^H(X,Y)Z-\Pi(Y,Z)\nabla_XN+\Pi(X,Z)\nabla_YN
    +\big(-(\nabla_X\Pi)(Y,Z)+(\nabla_Y\Pi)(X,Z)\big)N,
  \]
  where $(\nabla_X\Pi)(Y,Z)=D_X\Pi(Y,Z)-\Pi(\nabla^H_XY,Z)-\Pi(Y,\nabla^H_XZ)$ is the
  covariant derivative of $\Pi$ along $H$. Pairing with $g(-,W)$ for $W$ tangent
  to $H$, using $g(\nabla_XN,W)=\Pi(X,W)$ (\cref{def:pet-ch3-second-fundamental-form})
  and $g(N,W)=0$, gives
  \[
    g(R(X,Y)Z,W)=g_H(R^H(X,Y)Z,W)-\Pi(Y,Z)\Pi(X,W)+\Pi(X,Z)\Pi(Y,W).
  \]
\end{proof}

\begin{theorem}[Theorem 3.2.5 — the normal (Codazzi--Mainardi) curvature equation]
  \label{thm:pet-ch3-normal-curvature-equation}
  \lean{PetersenLib.normalCurvatureEquation}
  \leanok
  \uses{thm:pet-ch3-tangential-curvature-equation}
  For $X,Y,Z$ tangent to a hypersurface $H$,
  \[
    g(R(X,Y)Z,N)=-(\nabla_X\Pi)(Y,Z)+(\nabla_Y\Pi)(X,Z).
  \]
  \source{petersen:page-0111}
\end{theorem}
\begin{proof}
  \uses{thm:pet-ch3-normal-curvature-equation}
  Pair the expansion of $R(X,Y)Z$ obtained in the proof of
  \cref{thm:pet-ch3-tangential-curvature-equation} with $g(-,N)$ instead of
  $g(-,W)$. Since $R^H(X,Y)Z\in TH$, $g(R^H(X,Y)Z,N)=0$; since
  $g(\nabla_XN,N)=g(\nabla_YN,N)=0$ (\cref{def:pet-ch3-second-fundamental-form}),
  the $\Pi(Y,Z)\nabla_XN$ and $\Pi(X,Z)\nabla_YN$ terms drop out; and $g(N,N)=1$
  picks out the coefficient of $N$, giving
  $g(R(X,Y)Z,N)=-(\nabla_X\Pi)(Y,Z)+(\nabla_Y\Pi)(X,Z)$.
\end{proof}

\begin{remark}[curvature equations by dimension]
  \label{rem:pet-ch3-curvature-equations-by-dimension}
  \lean{PetersenLib.curvatureEquationsByDimension}
  \leanok
  \uses{thm:pet-ch3-radial-curvature-equation,thm:pet-ch3-tangential-curvature-equation,thm:pet-ch3-normal-curvature-equation}
  Together the three fundamental equations compute curvature by induction on
  dimension: knowing curvature on $H$ and the shape operator $S$ determines
  curvature on $M$ at points of $H$. If $\dim M=1$ then $\dim H=0$, matching
  $R\equiv0$ on all $1$-dimensional spaces. If $\dim M=2$ then $\dim H=1$, so
  $R^H\equiv0$ and $X,Y,Z$ are proportional; only the radial curvature equation
  is relevant. If $\dim M=3$ then $\dim H=2$: writing an orthonormal frame
  $\{X,Y,N\}$ with $X\perp Y$, the full curvature tensor depends on
  $g(R(X,N)N,Y)$, $g(R(X,N)N,X)$, $g(R(Y,N)N,Y)$ (radial equation),
  $g(R(X,Y)Y,X)$ (tangential equation), and $g(R(X,Y)Y,N)$, $g(R(Y,X)X,N)$
  (normal equation).
  \source{petersen:page-0112}
\end{remark}

\begin{corollary}[Gauss's Theorema Egregium]
  \label{cor:pet-ch3-gauss-egregium}
  \lean{PetersenLib.gaussTheoremaEgregium}
  \leanok
  \uses{thm:pet-ch3-tangential-curvature-equation,ex:pet-ch3-euclidean-flat}
  When $M^3=\mathbb R^3$ (so $R\equiv0$) and $E_1,E_2$ is an orthonormal basis of
  $T_pH$ for a surface $H\subset\mathbb R^3$,
  \[
    \sec(T_pH)=R^H(E_1,E_2,E_2,E_1)=\Pi(E_1,E_1)\Pi(E_2,E_2)-\Pi(E_1,E_2)^2=\det[\Pi],
  \]
  i.e.\ the extrinsic quantity $\det[\Pi]$ equals the intrinsic sectional
  curvature $\sec(T_pH)$.
  \source{petersen:page-0112}
\end{corollary}
\begin{proof}
  \uses{cor:pet-ch3-gauss-egregium}
  Apply \cref{thm:pet-ch3-tangential-curvature-equation} with $X=W=E_1$,
  $Y=Z=E_2$ and $R\equiv0$ (\cref{ex:pet-ch3-euclidean-flat}):
  \[
    0=g_H(R^H(E_1,E_2)E_2,E_1)-\Pi(E_1,E_1)\Pi(E_2,E_2)+\Pi(E_1,E_2)^2,
  \]
  and since $E_1,E_2$ are orthonormal, $\sec(T_pH)=g_H(R^H(E_1,E_2)E_2,E_1)$.
\end{proof}

\subsection{Distance Functions}

\begin{definition}[distance function]
  \label{def:pet-ch3-distance-function}
  \lean{PetersenLib.IsDistanceFunction}
  \leanok
  For $O\subset(M,g)$ open, $r:O\to\mathbb R$ is a \emph{distance function} if
  $|\nabla r|\equiv1$ on $O$, i.e.\ a solution of the \emph{Hamilton--Jacobi} (or
  \emph{eikonal}) equation $|\nabla r|^2=1$.
  \source{petersen:page-0113}
\end{definition}

\begin{example}[Example 3.2.6]
  \label{ex:pet-ch3-distance-to-points}
  \lean{PetersenLib.distanceFunction_points}
  \leanok
  \uses{def:pet-ch3-distance-function}
  On $(\mathbb R^n,g_{\mathrm{can}})$, $r(x)=|x-y|$ is smooth on $\mathbb
  R^n\setminus\{y\}$ with $|\nabla r|\equiv1$. For two points $y,z$,
  $r(x)=|x\{y,z\}|=\min\{|x-y|,|x-z|\}$ is smooth away from $\{y,z\}$ and the
  hyperplane equidistant from $y,z$.
  \source{petersen:page-0113}
\end{example}

\begin{example}[Example 3.2.7]
  \label{ex:pet-ch3-distance-to-submanifold}
  \lean{PetersenLib.distanceFunction_submanifold}
  \leanok
  % Lean: the flat / totally-geodesic case is formalized in full
  % (Ch03/DistanceFunctions.lean, distanceFunction_submanifold): for a linear
  % subspace H ⊆ ℝⁿ, r(x)=‖proj_{H⊥}x‖=‖x−proj_H x‖ is a distance function on
  % ℝⁿ∖H — smooth there and |∇r|≡1, with ∇r=proj_{H⊥}x/‖proj_{H⊥}x‖ the unit
  % normal (self-adjoint-idempotence of the orthogonal projection makes it a
  % unit vector).  The general curved-submanifold case r(x)=|xH| additionally
  % needs the tubular-neighbourhood / unique-nearest-point theorem to produce
  % the normal coordinates x=tN+y, which is not formalized.
  \uses{def:pet-ch3-distance-function}
  For a submanifold $H\subset\mathbb R^n$, $r(x)=|xH|=\inf\{|x-y|\mid y\in H\}$ is
  a distance function on some open $O$. When $H$ is an orientable hypersurface
  with unit normal $N$, coordinatize $\mathbb R^n$ by $x=tN+y$, $t\in\mathbb
  R,y\in H$; there is $\varepsilon:H\to(0,\infty)$ so that
  $O=\{tN+y\mid y\in H,|t|<\varepsilon(y)\}$ has unique coordinates $(t,y)$, and
  $r(x)=t$ (the \emph{signed distance}) or $f(x)=|xH|=|t|$ are distance functions
  on $O$, resp.\ $O\setminus H$. More generally, on $I\times H$ with a metric
  $dr^2+g_r$ ($g_r$ depending on $r$), the projection $I\times H\to I$ is a
  distance function; special cases include rotationally symmetric metrics,
  doubly warped products, and submersion metrics on $I\times S^{2n-1}$.
  \source{petersen:page-0113}
\end{example}

\begin{lemma}[Lemma 3.2.8]
  \label{lem:pet-ch3-distance-function-submersion}
  \lean{PetersenLib.distanceFunction_iff_riemannianSubmersion}
  \leanok
  \uses{def:pet-ch3-distance-function}
  For $O$ open in a Riemannian manifold, $r:O\to I\subset\mathbb R$ is a distance
  function if and only if it is a Riemannian submersion.
  \source{petersen:page-0114}
\end{lemma}
\begin{proof}
  \uses{lem:pet-ch3-distance-function-submersion}
  \leanok
  In general $Dr(v)=dr(v)\partial_t=g(\nabla r,v)\partial_t$, so $v\perp\ker Dr$
  iff $v$ is proportional to $\nabla r$. For $v=\alpha\nabla r$,
  $Dr(v)=\alpha\,Dr(\nabla r)=\alpha g(\nabla r,\nabla r)\partial_t$. Since
  $\partial_t$ has length $1$ in $I$, $|v|=|\alpha||\nabla r|$ and
  $|Dr(v)|=|\alpha||\nabla r|^2$. Thus $Dr$ is a linear isometry on
  $(\ker Dr)^\perp$ iff $|\nabla r|=1$, i.e.\ $r$ is a Riemannian submersion iff
  $|\nabla r|\equiv1$.
\end{proof}

\begin{definition}[distance-function notation: $\partial_r,O_r,g_r,S$]
  \label{def:pet-ch3-distance-function-notation}
  \lean{PetersenLib.distanceFunctionNotation}
  \leanok
  \uses{lem:pet-ch3-distance-function-submersion,def:pet-ch3-second-fundamental-form}
  Fix a distance function $r:O\to\mathbb R$. Write $\partial_r=\nabla r$; the
  level sets are $O_r=\{x\in O\mid r(x)=r\}$ with induced metric $g_r$, and
  $\nabla',R'$ denote the connection and curvature of $(O_r,g_r)$. Since
  $|\nabla r|=1$, $\operatorname{Hess}r=\Pi$, and $S$ denotes the common
  $(1,1)$-tensor for both: $S=\nabla\partial_r$ measures how the induced metric
  on $O_r$ changes, by measuring how the unit normal to $O_r$ changes.
  \source{petersen:page-0114}
\end{definition}

\begin{example}[Example 3.2.9]
  \label{ex:pet-ch3-zero-shape-operator-hyperplane}
  \lean{PetersenLib.zeroShapeOperator_iff_hyperplane}
  \leanok
  \uses{def:pet-ch3-distance-function-notation}
  Let $H\subset\mathbb R^n$ be an orientable hypersurface with unit normal $N$
  and shape operator $S(v)=\nabla_vN$. If $S=0$ on $H$, $N$ is constant on $H$
  and $H$ is open in the hyperplane $\{x+p\mid x\cdot N_p=0\}$ for fixed $p\in
  H$. Explicitly, for the isometric immersion
  $(x^1,\dots,x^{n-1})\mapsto(c(x^1),x^2,\dots,x^{n-1})$ of a unit-speed curve $c$
  into $\mathbb R^n$, the unit normal is $N=(-\dot c^2(x^1),\dot c^1(x^1),0,\dots,0)$,
  and $S\equiv0$ iff $\dot c^1\equiv\dot c^2\equiv0$ iff $c$ is a straight line
  iff $H$ is open in a hyperplane.
  \source{petersen:page-0114,petersen:page-0115}
\end{example}
\begin{proof}
  \uses{ex:pet-ch3-zero-shape-operator-hyperplane}
  Since $N$ has unit length in Cartesian coordinates,
  \[
    \nabla N=-d(\dot c^2)\partial_1+d(\dot c^1)\partial_2
    =-\ddot c^2dx^1\partial_1+\ddot c^1dx^1\partial_2
    =(-\ddot c^2\partial_1+\ddot c^1\partial_2)dx^1,
  \]
  so $S\equiv0$ iff $\ddot c^1\equiv\ddot c^2\equiv0$, i.e.\ $c$ is a straight
  line, i.e.\ the image of the immersion is (an open subset of) a hyperplane.
  Conversely, if $S=0$ on a general orientable hypersurface $H\subset\mathbb
  R^n$, then $N$ is a parallel, hence constant, vector field on $H$, and $H$ is
  tangent everywhere to the fixed hyperplane $x\cdot N_p=\text{const}$ through
  any of its points $p$, hence is an open subset of that hyperplane.
\end{proof}

\begin{remark}[intrinsic versus extrinsic geometry]
  \label{rem:pet-ch3-intrinsic-extrinsic}
  \lean{PetersenLib.intrinsicExtrinsicGeometry}
  \leanok
  \uses{ex:pet-ch3-zero-shape-operator-hyperplane}
  Intrinsic geometry of $(M,g)$ is everything computable without reference to an
  ambient isometric immersion; extrinsic geometry studies how an isometric
  immersion $(M,g)\to(\bar M,\bar g)$ bends $(M,g)$ inside $(\bar M,\bar g)$. The
  curvature tensor on $(M,g)$ measures intrinsic bending, while the shape
  operator measures extrinsic bending.
  \source{petersen:page-0115}
\end{remark}

\subsection{The Curvature Equations for Distance Functions}

\begin{corollary}[Corollary 3.2.10]
  \label{cor:pet-ch3-radial-equation-distance-function}
  \lean{PetersenLib.radialEquation_distanceFunction}
  \leanok
  \uses{def:pet-ch3-distance-function-notation,thm:pet-ch3-radial-curvature-equation,prop:pet-ch3-normal-hessian-relation}
  For a distance function $r:O\to\mathbb R$, $\nabla_{\partial_r}\partial_r=0$ and
  $\nabla_{\partial_r}S+S^2=-R_{\partial_r}$, where $R_{\partial_r}(X)=R(X,\partial_r)\partial_r$.
  \source{petersen:page-0115}
\end{corollary}
\begin{proof}
  \uses{cor:pet-ch3-radial-equation-distance-function}
  \leanok
  The first fact is part (3) of \cref{prop:pet-ch3-normal-hessian-relation}
  applied to $f=r$: $\operatorname{Hess}r(\nabla r,X)=\tfrac12D_X|\nabla r|^2=0$
  since $|\nabla r|\equiv1$, i.e.\ $g(\nabla_{\partial_r}\partial_r,X)=0$ for all
  $X$, so $\nabla_{\partial_r}\partial_r=0$. The second is the first formula of
  \cref{thm:pet-ch3-radial-curvature-equation} for $f=r$: since
  $\nabla_{\partial_r}\partial_r=0$, the term
  $\nabla_X(S(\nabla f))=\nabla_X(S(\partial_r))=\nabla_X(\nabla_{\partial_r}\partial_r)=0$
  drops out, leaving $(\nabla_{\partial_r}S)(X)+S^2(X)=-R(X,\partial_r)\partial_r$.
\end{proof}

\begin{proposition}[Proposition 3.2.11]
  \label{prop:pet-ch3-distance-function-curvature-equations}
  \lean{PetersenLib.distanceFunction_curvatureEquations}
  \leanok
  \uses{cor:pet-ch3-radial-equation-distance-function,thm:pet-ch3-radial-curvature-equation}
  For a distance function $r:(O,g)\to\mathbb R$ with $\partial_r=\nabla r$:
  \begin{enumerate}
    \item[(1)] $\mathcal L_{\partial_r}g=2\operatorname{Hess}r$;
    \item[(2)] $(\nabla_{\partial_r}\operatorname{Hess}r)(X,Y)+\operatorname{Hess}^2r(X,Y)=-R(X,\partial_r,\partial_r,Y)$;
    \item[(3)] $(\mathcal L_{\partial_r}\operatorname{Hess}r)(X,Y)-\operatorname{Hess}^2r(X,Y)=-R(X,\partial_r,\partial_r,Y)$.
  \end{enumerate}
  \source{petersen:page-0115,petersen:page-0116}
\end{proposition}
\begin{proof}
  \uses{prop:pet-ch3-distance-function-curvature-equations}
  \leanok
  (1) is the standard identity $\operatorname{Hess}f=\tfrac12\mathcal
  L_{\nabla f}g$ applied to $f=r$. (2) and (3) follow directly from
  \cref{thm:pet-ch3-radial-curvature-equation} (second and third formulas)
  applied to $f=r$, using $|\nabla r|\equiv1$ so that
  $\operatorname{Hess}(\tfrac12|\nabla r|^2)=\operatorname{Hess}(\tfrac12)=0$.
\end{proof}

\subsection{Jacobi Fields}

\begin{definition}[Jacobi field]
  \label{def:pet-ch3-jacobi-field}
  \lean{PetersenLib.IsJacobiField}
  \leanok
  \uses{def:pet-ch3-distance-function-notation}
  For a smooth distance function $r$, a \emph{Jacobi field} is a smooth vector
  field $J$ independent of $r$, i.e.\ satisfying the \emph{Jacobi equation}
  $\mathcal L_{\partial_r}J=0$.
  \source{petersen:page-0116}
\end{definition}

\begin{remark}[construction of Jacobi fields]
  \label{rem:pet-ch3-jacobi-field-construction}
  \lean{PetersenLib.jacobiField_construction}
  \leanok
  \uses{def:pet-ch3-jacobi-field}
  The Jacobi equation is a first-order linear PDE solvable by characteristics: in
  local coordinates $(r,x^2,\dots,x^n)$, $J=J^r\partial_r+J^i\partial_i$, and
  $\mathcal L_{\partial_r}J=\partial_r(J^r)\partial_r+\partial_r(J^i)\partial_i$, so the Jacobi
  equation is equivalent to $J^r,J^i$ being independent of $r$; hence a Jacobi
  field is determined by its values on a hypersurface $H\subset M$ transverse to
  $\partial_r$ where $(x^2,\dots,x^n)|_H$ is a coordinate system. Coordinate
  vector fields are themselves Jacobi fields.
  \source{petersen:page-0116}
\end{remark}

\begin{lemma}[first-order form of the Jacobi equation]
  \label{lem:pet-ch3-jacobi-field-equivalent-equation}
  \lean{PetersenLib.jacobiField_equivalentEquation}
  \leanok
  \uses{def:pet-ch3-jacobi-field}
  $\mathcal L_{\partial_r}J=0$ is equivalent to $\nabla_{\partial_r}J=S(J)$, and
  consequently $\operatorname{Hess}r(J,J)=g(\nabla_{\partial_r}J,J)=\tfrac12\partial_rg(J,J)$.
  \source{petersen:page-0116}
\end{lemma}
\begin{proof}
  \uses{lem:pet-ch3-jacobi-field-equivalent-equation}
  \leanok
  For vector fields, $\mathcal L_{\partial_r}J=[\partial_r,J]=\nabla_{\partial_r}J-\nabla_J\partial_r$
  (torsion-freeness), and $\nabla_J\partial_r=S(J)$ by definition of $S$. Hence
  $\mathcal L_{\partial_r}J=0$ iff $\nabla_{\partial_r}J=S(J)$. Then $g(\nabla_{\partial_r}J,J)=g(S(J),J)=\operatorname{Hess}r(J,J)$,
  while $g(\nabla_{\partial_r}J,J)=\tfrac12D_{\partial_r}g(J,J)=\tfrac12\partial_rg(J,J)$.
\end{proof}

\begin{proposition}[second-order Jacobi equation]
  \label{prop:pet-ch3-jacobi-second-order-equation}
  \lean{PetersenLib.jacobiField_secondOrderEquation}
  \leanok
  \uses{lem:pet-ch3-jacobi-field-equivalent-equation,cor:pet-ch3-radial-equation-distance-function,def:pet-ch3-directional-curvature-operator}
  A Jacobi field $J$ satisfies $\nabla_{\partial_r}\nabla_{\partial_r}J=-R(J,\partial_r)\partial_r$.
  \source{petersen:page-0117}
\end{proposition}
\begin{proof}
  \uses{prop:pet-ch3-jacobi-second-order-equation}
  \leanok
  Using the Jacobi equation $\nabla_{\partial_r}J=S(J)$ and the product rule,
  \[
    \nabla_{\partial_r}\nabla_{\partial_r}J=\nabla_{\partial_r}(S(J))
    =(\nabla_{\partial_r}S)(J)+S(\nabla_{\partial_r}J)=(\nabla_{\partial_r}S)(J)+S^2(J).
  \]
  By \cref{cor:pet-ch3-radial-equation-distance-function},
  $(\nabla_{\partial_r}S)(J)+S^2(J)=-R_{\partial_r}(J)=-R(J,\partial_r)\partial_r$
  (\cref{def:pet-ch3-directional-curvature-operator}).
\end{proof}

\begin{proposition}[fundamental equations on Jacobi fields]
  \label{prop:pet-ch3-jacobi-field-curvature-equations}
  \lean{PetersenLib.jacobiField_curvatureEquations}
  \leanok
  \uses{def:pet-ch3-jacobi-field,prop:pet-ch3-distance-function-curvature-equations}
  For Jacobi fields $J_1,J_2$, equations (1) and (3) of
  \cref{prop:pet-ch3-distance-function-curvature-equations} become
  \[
    \partial_rg(J_1,J_2)=2\operatorname{Hess}r(J_1,J_2),\qquad
    \partial_r\big(\operatorname{Hess}r(J_1,J_2)\big)-\operatorname{Hess}^2r(J_1,J_2)=-R(J_1,\partial_r,\partial_r,J_2).
  \]
  \source{petersen:page-0117}
\end{proposition}
\begin{proof}
  \uses{prop:pet-ch3-jacobi-field-curvature-equations}
  \leanok
  For a function $\phi=g(J_1,J_2)$ on $M$, $\mathcal
  L_{\partial_r}\phi=D_{\partial_r}\phi=\partial_r\phi$ always; also
  $\mathcal L_{\partial_r}(g(J_1,J_2))=(\mathcal L_{\partial_r}g)(J_1,J_2)+g(\mathcal
  L_{\partial_r}J_1,J_2)+g(J_1,\mathcal L_{\partial_r}J_2)=(\mathcal
  L_{\partial_r}g)(J_1,J_2)$ since $J_1,J_2$ are Jacobi fields. Combined with (1)
  of \cref{prop:pet-ch3-distance-function-curvature-equations} this gives the
  first displayed equation. Likewise
  $\mathcal L_{\partial_r}(\operatorname{Hess}r(J_1,J_2))=(\mathcal
  L_{\partial_r}\operatorname{Hess}r)(J_1,J_2)=\partial_r(\operatorname{Hess}r(J_1,J_2))$,
  and substituting into (3) gives the second equation.
\end{proof}

\begin{remark}[reduction on a product neighborhood]
  \label{rem:pet-ch3-jacobi-field-product-neighborhood}
  \lean{PetersenLib.jacobiField_productNeighborhood}
  \leanok
  \uses{prop:pet-ch3-jacobi-field-curvature-equations}
  On a product neighborhood $\Omega=(a,b)\times H\subset M$ with $r(t,z)=t$, the
  metric is $g=dr^2+g_r$, a one-parameter family of metrics on $H$; a vector
  field on $H$ extends uniquely to a Jacobi field on $\Omega$. Since
  $\operatorname{Hess}r(\partial_r,J)=g(\nabla_{\partial_r}\partial_r,J)=0$ and
  $g_r(\partial_r,J)=0$, restricting to $H$ gives $\partial_rg=\partial_rg_r=2\operatorname{Hess}r$,
  and the fundamental equations become, on $H$,
  \[
    \partial_rg_r=2\operatorname{Hess}r,\qquad
    \partial_r\operatorname{Hess}r-\operatorname{Hess}^2r=-R(\cdot,\partial_r,\partial_r,\cdot).
  \]
  \source{petersen:page-0117}
\end{remark}

\begin{lemma}[the radial curvature term via sectional curvature]
  \label{lem:pet-ch3-jacobi-field-sectional-curvature-formula}
  \lean{PetersenLib.jacobiField_sectionalCurvatureFormula}
  \leanok
  \uses{rem:pet-ch3-jacobi-field-product-neighborhood,def:pet-ch3-sectional-curvature}
  For $X$ tangent to $H$ (so $g(X,\partial_r)=0$),
  $R(X,\partial_r,\partial_r,X)=\sec(X,\partial_r)g_r(X,X)$; hence for a Jacobi
  field $J$ on $H$,
  \[
    \partial_r\big(\operatorname{Hess}r(J,J)\big)-\operatorname{Hess}^2r(J,J)=-\sec(J,\partial_r)g_r(J,J).
  \]
  This still couples the Hessian equation to the metric $g_r$.
  \source{petersen:page-0118}
\end{lemma}
\begin{proof}
  \uses{lem:pet-ch3-jacobi-field-sectional-curvature-formula}
  \leanok
  Since $g(X,\partial_r)=0$ and $g(\partial_r,\partial_r)=1$,
  \[
    R(X,\partial_r,\partial_r,X)=\sec(X,\partial_r)\big(g(X,X)g(\partial_r,\partial_r)-g(X,\partial_r)^2\big)
    =\sec(X,\partial_r)g(X,X)=\sec(X,\partial_r)g_r(X,X),
  \]
  by \cref{def:pet-ch3-sectional-curvature}. Substituting $X=J$ into the second
  equation of \cref{rem:pet-ch3-jacobi-field-product-neighborhood} evaluated on
  $(J,J)$ gives the displayed identity.
\end{proof}

\begin{example}[Example 3.2.12]
  \label{ex:pet-ch3-two-dim-sectional-curvature-formula}
  \lean{PetersenLib.twoDimensional_sectionalCurvature_formula}
  \leanok
  \uses{lem:pet-ch3-jacobi-field-sectional-curvature-formula}
  If $\dim M=2$, write $g=dr^2+\rho^2(r,\theta)d\theta^2$ with $\theta$
  coordinatizing the level sets of $r$. Then $\partial_\theta$ is a Jacobi field
  of length $\rho$, and $\sec(T_pM)=-\partial_r^2\rho/\rho$.
  \source{petersen:page-0118,petersen:page-0119}
\end{example}
\begin{proof}
  \uses{ex:pet-ch3-two-dim-sectional-curvature-formula}
  \leanok
  Since $\partial_\theta$ is a Jacobi field (a coordinate vector field, by
  \cref{rem:pet-ch3-jacobi-field-construction}), the first formula of
  \cref{prop:pet-ch3-jacobi-field-curvature-equations} gives
  $2\operatorname{Hess}r(\partial_\theta,\partial_\theta)=\partial_r\rho^2=2\rho\partial_r\rho$, so
  $\operatorname{Hess}r(\partial_\theta,\partial_\theta)=\rho\partial_r\rho$. Since $[\partial_r,\partial_\theta]=0$
  and $S$ is self-adjoint with $S(\partial_r)=0$,
  $S(\partial_\theta)=\dfrac{\partial_r\rho}{\rho}\partial_\theta$. By
  \cref{lem:pet-ch3-jacobi-field-sectional-curvature-formula} applied to
  $J=\partial_\theta$ (so $g_r(J,J)=\rho^2$),
  \begin{align*}
    -\sec(\partial_\theta,\partial_r)\rho^2
    &=\partial_r\big(\operatorname{Hess}r(\partial_\theta,\partial_\theta)\big)-\operatorname{Hess}^2r(\partial_\theta,\partial_\theta)\\
    &=\partial_r(\rho\partial_r\rho)-(\partial_r\rho)^2=\rho\partial_r^2\rho,
  \end{align*}
  giving $\sec(T_pM)=\sec(\partial_\theta,\partial_r)=-\partial_r^2\rho/\rho$.
\end{proof}

\subsection{Parallel Fields}

\begin{definition}[parallel field]
  \label{def:pet-ch3-parallel-field}
  \lean{PetersenLib.IsParallelField}
  \leanok
  \uses{def:pet-ch3-distance-function-notation}
  For a smooth distance function $r$, a \emph{parallel field} is a vector field
  $X$ with $\nabla_{\partial_r}X=0$. Parallel fields are almost never Jacobi
  fields.
  \source{petersen:page-0119}
\end{definition}

\begin{proposition}[curvature equation on parallel fields]
  \label{prop:pet-ch3-parallel-field-curvature-equation}
  \lean{PetersenLib.parallelField_curvatureEquation}
  \leanok
  \uses{def:pet-ch3-parallel-field,prop:pet-ch3-distance-function-curvature-equations,lem:pet-ch3-jacobi-field-sectional-curvature-formula}
  For parallel fields $X,Y$, $\partial_rg(X,Y)=0$; equation (2) of
  \cref{prop:pet-ch3-distance-function-curvature-equations} becomes
  $\partial_r(\operatorname{Hess}r(X,Y))+\operatorname{Hess}^2r(X,Y)=-R(X,\partial_r,\partial_r,Y)$,
  and on a unit parallel field $X$ orthogonal to $\partial_r$,
  \[
    \partial_r\big(\operatorname{Hess}r(X,X)\big)+\operatorname{Hess}^2r(X,X)=-\sec(X,\partial_r),
  \]
  decoupled from the metric since $g_r(X,X)$ is then constant in $r$.
  \source{petersen:page-0119}
\end{proposition}
\begin{proof}
  \uses{prop:pet-ch3-parallel-field-curvature-equation}
  \leanok
  For parallel $X,Y$, $\partial_rg(X,Y)=g(\nabla_{\partial_r}X,Y)+g(X,\nabla_{\partial_r}Y)=0$.
  Since $X,Y$ are constant in $r$, they contribute no extra terms to (2) of
  \cref{prop:pet-ch3-distance-function-curvature-equations} beyond the covariant
  derivative form already displayed. Projecting $X$ onto $H$ and rescaling gives
  a unit parallel field orthogonal to $\partial_r$ with $g_r(X,X)\equiv1$
  constant; by \cref{lem:pet-ch3-jacobi-field-sectional-curvature-formula} (whose
  proof used only orthogonality to $\partial_r$, not the Jacobi equation),
  $R(X,\partial_r,\partial_r,X)=\sec(X,\partial_r)g_r(X,X)=\sec(X,\partial_r)$, giving the
  displayed decoupled equation.
\end{proof}

\subsection{Conjugate Points}

\begin{definition}[conjugate point, focal point]
  \label{def:pet-ch3-conjugate-focal-point}
  \lean{PetersenLib.ConjugatePoint}
  \leanok
  \uses{prop:pet-ch3-distance-function-curvature-equations,prop:pet-ch3-parallel-field-curvature-equation}
  Along an integral curve of $\partial_r$ with bounded curvature, if
  $\operatorname{Hess}r$ blows up, equation (2) forces it to be decreasing as $r$
  increases, so it can only tend to $-\infty$; by equation (1), the only possible
  degeneration along the curve is that the metric $g$ stops being positive
  definite. When this happens $r$ develops a \emph{conjugate} or \emph{focal
  point} along the curve. A \emph{conjugate point} occurs when
  $\operatorname{Hess}r$ becomes undefined while solving its differential
  equation; a \emph{focal point} occurs when integral curves of $\nabla r$ meet.
  The two frequently coincide but need not: there are metrics with conjugate
  points that are not focal points.
  \source{petersen:page-0120}
\end{definition}

\begin{remark}[conjugate versus focal points for an ellipse]
  \label{rem:pet-ch3-conjugate-vs-focal-example}
  \lean{PetersenLib.conjugateVsFocal_ellipse}
  \uses{def:pet-ch3-conjugate-focal-point}
  For the lower half of an ellipse, conjugate points occur along its evolute.
  For the entire ellipse, the focal set is the segment between the two focal
  points of the ellipse, since the normal lines from the top and bottom arcs
  intersect along this segment.
  \source{petersen:page-0120}
\end{remark}

\begin{remark}[external force versus internal reaction]
  \label{rem:pet-ch3-external-force-internal-reaction}
  \lean{PetersenLib.externalForceInternalReaction}
  \leanok
  \uses{def:pet-ch3-conjugate-focal-point,prop:pet-ch3-distance-function-curvature-equations}
  Rewriting equations (2) and (3) of
  \cref{prop:pet-ch3-distance-function-curvature-equations} as
  \[
    (\nabla_{\partial_r}\operatorname{Hess}r)(X,X)=-R(X,\partial_r,\partial_r,X)-\operatorname{Hess}^2r(X,X),\qquad
    (\mathcal L_{\partial_r}\operatorname{Hess}r)(X,X)=-R(X,\partial_r,\partial_r,X)+\operatorname{Hess}^2r(X,X),
  \]
  the curvature term is a fixed \emph{external force}, while
  $\operatorname{Hess}^2r$ (of fixed sign) is an \emph{internal reaction} that
  forces $\operatorname{Hess}r$ to blow up or collapse in finite time. If
  $\sec\le0$, $\mathcal L_{\partial_r}\operatorname{Hess}r\ge0$: positive
  $\operatorname{Hess}r$ stays positive. If $\sec\ge0$,
  $\nabla_{\partial_r}\operatorname{Hess}r\le0$: nonpositive
  $\operatorname{Hess}r$ stays nonpositive.
  \source{petersen:page-0120}
\end{remark}

\section{Further Study}

Two classical references are singled out for further reading on the tensorial
aspects of Riemannian and pseudo-Riemannian geometry: an older but thorough
treatment recommended to anyone learning the subject, and an authoritative
multi-volume guide (particularly its second volume; the first is best used as a
reference for foundational material).

\section{Exercises}

\begin{remark}[Exercise 3.4.1]
  \label{rem:pet-ch3-ex-1}
  \lean{PetersenLib.exercise3_4_1}
  \leanok
  % Lean: self-contained coordinate proof (Ch03/EuclideanImmersionCurvature.lean),
  % decoupled from the manifold machinery.  The induced metric g_{ij}=⟨∂_iu,∂_ju⟩,
  % its inverse, the Christoffel symbols Γ_{k,ij}=⟨∂²_{ij}u,∂_ku⟩ (= the usual
  % ½(∂g+∂g−∂g)), and the induced connection ∇^M are built as plain fderiv/matrix
  % analysis; the second fundamental form II_{ij}=∂²_{ij}u−∇^M_{∂_i}∂_j is the
  % normal part.  `gaussEquation` proves the Gauss / Theorema Egregium identity
  % R_{ijkl}=⟨II_{jk},II_{il}⟩−⟨II_{ik},II_{jl}⟩ — the third derivatives cancel by
  % Clairaut symmetry, and the projection is used only *through* the pairing
  % ⟨II,∂_lu⟩=0, so the inverse Gram matrix is never differentiated.  The node
  % itself (`exercise3_4_1`) is the resulting **jet invariance**: two smooth
  % parametrizations with equal first and second partials at a point induce the
  % same mixed curvature R^l_{ijk}=∑_p g^{lp}R_{ijkp} there.
  \uses{def:pet-ch3-curvature-tensor}
  Let $M$ be an $n$-dimensional submanifold of $\mathbb R^{n+m}$ with the induced
  metric, given locally by a parametrization $u^s(x^1,\dots,x^n)$,
  $s=1,\dots,n+m$. Show that in these coordinates $R^l_{ijk}$ depends only on the
  first and second partials of $u^s$.
  \source{petersen:page-0121}
\end{remark}

\begin{remark}[Exercise 3.4.2]
  \label{rem:pet-ch3-ex-2}
  \lean{PetersenLib.exercise3_4_2}
  \leanok
  \uses{def:pet-ch3-hessian-square-notation}
  For $f:(M,g)\to\mathbb R$ smooth on a connected Riemannian manifold, consider:
  (1) $|\nabla f|$ is constant; (2) $\nabla_{\nabla f}\nabla f=0$; (3) $|\nabla
  f|$ is constant on the level sets of $f$. Show $(1)\Leftrightarrow(2)\Rightarrow(3)$,
  and give an example showing the last implication is not an equivalence.
  \source{petersen:page-0121}
\end{remark}

\begin{remark}[Exercise 3.4.3]
  \label{rem:pet-ch3-ex-3}
  \lean{PetersenLib.exercise3_4_3}
  \leanok
  \uses{def:pet-ch3-hessian-square-notation,thm:pet-ch3-radial-curvature-equation}
  For $f$ with $S(X)=\nabla_X\nabla f$, show that
  $\mathcal L_{\nabla f}S=\nabla_{\nabla f}S$ and
  $\mathcal L_{\nabla f}S+S^2-\nabla_X(S(\nabla f))=-R_{\nabla f}$. Reconcile this
  with \cref{thm:pet-ch3-radial-curvature-equation} for the $(0,2)$-version of
  the Hessian.
  \source{petersen:page-0121}
\end{remark}

\begin{remark}[Exercise 3.4.4]
  \label{rem:pet-ch3-ex-4}
  \lean{PetersenLib.exercise3_4_4}
  \leanok
  \uses{thm:pet-ch3-tangential-curvature-equation,thm:pet-ch3-normal-curvature-equation}
  For a distance function $r=f:M\to\mathbb R$, show the tangential and mixed
  curvature equations can be written as
  \[
    (R(X,Y)Z)^\top=R_H(X,Y)Z-(S(X)\wedge S(Y))(Z),\qquad
    g(R(X,Y)Z,N)=-g((\nabla_XS)(Y),Z)+g((\nabla_YS)(X),Z),
  \]
  and $R(X,Y)N=(d^\nabla S)(X,Y)$.
  \source{petersen:page-0121,petersen:page-0122}
\end{remark}

\begin{remark}[Exercise 3.4.5]
  \label{rem:pet-ch3-ex-5}
  \lean{PetersenLib.exercise3_4_5}
  \leanok
  \uses{prop:pet-ch3-curvature-symmetries}
  Prove both Bianchi identities at a point $p\in M$ using a coordinate system
  where $\nabla_{\partial_i}\partial_j=0$ at $p$.
  \source{petersen:page-0122}
\end{remark}

\begin{remark}[Exercise 3.4.6]
  \label{rem:pet-ch3-ex-6}
  \lean{PetersenLib.exercise3_4_6}
  \leanok
  \uses{def:pet-ch3-constant-curvature}
  Show that a Riemannian manifold with constant curvature has parallel curvature
  tensor.
  \source{petersen:page-0122}
\end{remark}

\begin{remark}[Exercise 3.4.7]
  \label{rem:pet-ch3-ex-7}
  \lean{PetersenLib.exercise3_4_7}
  \leanok
  \uses{def:pet-ch3-ricci-curvature,def:pet-ch3-scalar-curvature}
  Show that a Riemannian manifold with parallel Ricci tensor has constant scalar
  curvature. (The converse fails, and a metric with parallel curvature tensor
  need not be Einstein.)
  \source{petersen:page-0122}
\end{remark}

\begin{remark}[Exercise 3.4.8]
  \label{rem:pet-ch3-ex-8}
  \lean{PetersenLib.exercise3_4_8}
  \leanok
  \uses{prop:pet-ch3-contracted-bianchi,notation:pet-ch3-divergence-adjoint}
  With $R$ the $(0,4)$-curvature tensor and $\operatorname{Ric}$ the
  $(0,2)$-Ricci tensor, show
  $(\nabla^*R)(Z,X,Y)=(\nabla_X\operatorname{Ric})(Y,Z)-(\nabla_Y\operatorname{Ric})(X,Z)$.
  Conclude $\nabla^*R=0$ if $\nabla\operatorname{Ric}=0$, and that $\nabla^*R=0$
  iff the $(1,1)$-Ricci tensor satisfies $(\nabla_X\operatorname{Ric})(Y)=(\nabla_Y\operatorname{Ric})(X)$
  for all $X,Y$.
  \source{petersen:page-0122}
\end{remark}

\begin{remark}[Exercise 3.4.9]
  \label{rem:pet-ch3-ex-9}
  \lean{PetersenLib.exercise3_4_9}
  \uses{def:pet-ch3-sectional-curvature}
  For Riemannian manifolds $(M,g_M),(N,g_N)$ with product metric
  $(M\times N,g_M+g_N)$, and $X$ a vector field on $M$, $Y$ on $N$ regarded on
  $M\times N$, show $\nabla_XY=0$; conclude $\sec(X,Y)=0$, so product metrics
  always have vanishing curvatures.
  \source{petersen:page-0122}
\end{remark}

\begin{remark}[Exercise 3.4.10]
  \label{rem:pet-ch3-ex-10}
  \lean{PetersenLib.exercise3_4_10}
  \leanok
  \uses{def:pet-ch3-constant-curvature}
  Show that $(M,g)$ has constant curvature at $p$ iff $R(v,w)z=0$ for all pairwise
  orthogonal $v,w,z\in T_pM$. Hint: a symmetric bilinear form $B$ on an inner
  product space with $B(v,w)=0$ whenever $v\perp w$ is a multiple of the inner
  product.
  \source{petersen:page-0122}
\end{remark}

\begin{remark}[Exercise 3.4.11]
  \label{rem:pet-ch3-ex-11}
  \lean{PetersenLib.exercise3_4_11}
  \leanok
  % Lean: the tangential/normal splitting of the ambient tangent bundle is
  % encoded, exactly as in exercise3_4_14, by the smooth field of orthogonal
  % projections P onto the normal bundle; ∇^M = shapeTangent = (∇·)^⊤ is the
  % induced connection, T = normalCov = (∇·)^⊥ the (vector-valued) second
  % fundamental form, ∇^⊥ = normalCov the normal connection.  Since ∇_v W splits
  % *definitionally* into (∇_v W)^⊤ + (∇_v W)^⊥, expanding each second covariant
  % derivative gives the four (⊤/⊥)×(⊤/⊥) blocks that assemble into R^M, the
  % shape terms, and the ∇^⊥T terms; the only non-formal inputs are additivity
  % of ∇ over a field split, torsion-freeness, and integrability P([X,Y])=0 of
  % the tangent distribution (automatic on a genuine submanifold).  The
  % statement is proved as a pure operator identity for any smooth X,Y,Z with
  % P([X,Y]_p)=0; no idempotence/self-adjointness of P is needed for it.
  \uses{thm:pet-ch3-tangential-curvature-equation}
  For $X,Y,Z$ tangent to $M$ (a submanifold with second fundamental form $T$ and
  normal connection $\nabla^\perp$), show
  \[
    \bar R(X,Y)Z=R^M(X,Y)Z+T_XT_YZ-T_YT_XZ+(\nabla^\perp_XT)_YZ-(\nabla^\perp_YT)_XZ,
  \]
  where $\bar R$ is the curvature of the ambient manifold and
  $(\nabla^\perp_XT)_YZ=\nabla^\perp_X(T_YZ)-T_{\nabla^M_XY}Z-T_Y\nabla^M_XZ$;
  the tangential parts recover the Gauss equations, the normal parts the
  Peterson--Codazzi--Mainardi equations. (The source prints $R^M$ on both sides;
  the left-hand curvature is the ambient one $\bar R$, not the induced $R^M$.)
  \source{petersen:page-0122}
\end{remark}

\begin{remark}[Exercise 3.4.12]
  \label{rem:pet-ch3-ex-12}
  \lean{PetersenLib.exercise3_4_12}
  \leanok
  \uses{def:pet-ch3-second-fundamental-form,def:pet-ch3-ricci-curvature}
  For $H^{n-1}\subset\mathbb R^n$, show
  $\operatorname{Ric}^H=\operatorname{tr}(\Pi)\cdot\Pi-\Pi^2$.
  \source{petersen:page-0123}
\end{remark}

\begin{remark}[Exercise 3.4.13]
  \label{rem:pet-ch3-ex-13}
  \lean{PetersenLib.exercise3_4_13}
  \leanok
  % Lean: part (2) — the property forces constant curvature (n>2); part (1) needs space-form models (not yet formalized).
  \uses{def:pet-ch3-second-fundamental-form,def:pet-ch3-constant-curvature,rem:pet-ch3-ex-10}
  A hypersurface is \emph{totally geodesic} if its second fundamental form
  vanishes. (1) Show that the space forms $S^n_k$ have the property that any
  tangent vector is normal to a totally geodesic hypersurface. (2) Show that a
  Riemannian $n$-manifold, $n>2$, with this property has constant curvature.
  Hint: show $R(X,Y)Z=0$ for pairwise orthogonal $X,Y,Z$ and use
  \cref{rem:pet-ch3-ex-10}.
  \source{petersen:page-0123}
\end{remark}

\begin{remark}[Exercise 3.4.14]
  \label{rem:pet-ch3-ex-14}
  \lean{PetersenLib.exercise3_4_14}
  \leanok
  % Lean: the general-codimension normal bundle is encoded abstractly by its
  % orthogonal projection P (idempotent, g-self-adjoint, smooth) onto the normal
  % bundle over an open U; normal fields are the sections fixed by P. From P we
  % build the normal connection ∇^⊥ = P∘∇, the tangential shape form T_vV = ∇_vV
  % − P(∇_vV), and the normal curvature R^⊥. exercise3_4_14 proves (1) the two
  % skew-symmetries (in X,Y and in V,W) and (2) the Ricci equations, for any
  % smooth orthogonal splitting (no integrability of the tangent distribution
  % needed). Tensoriality of R^⊥ (also part (1)) is exhibited by the Ricci
  % equations, which write R^⊥ = R^M − (shape-form difference), a difference of
  % genuine (0,4)-tensors; the substantive content — the skew-symmetries — is
  % formalised, full C^∞-linearity is not re-derived.
  \uses{def:pet-ch3-second-fundamental-form}
  For tangent fields $X,Y$ and normal fields $V,W$ of a submanifold, define the
  normal curvature $R^\perp(X,Y,V,W)$. (1) Show $R^\perp$ is tensorial and
  skew-symmetric in $X,Y$ and in $V,W$. (2) Show
  $R^M(X,Y,V,W)=R^\perp(X,Y,V,W)+g_M(T_XV,T_YW)-g_M(T_YV,T_XW)$ (the
  \emph{Ricci equations}).
  (The Lean formalisation encodes the tangent/normal splitting abstractly by the
  orthogonal projection onto the normal bundle over an open set, and proves the
  two skew-symmetries of part~(1) together with the Ricci equations of part~(2)
  for any such smooth orthogonal splitting; the tensoriality of $R^\perp$ is
  exhibited by the Ricci equations — which write $R^\perp$ as a difference of the
  ambient $(0,4)$-curvature and a shape-form tensor — rather than re-derived as a
  separate $C^\infty$-linearity statement.)
  \source{petersen:page-0123}
\end{remark}

\begin{remark}[Exercise 3.4.15]
  \label{rem:pet-ch3-ex-15}
  \lean{PetersenLib.exercise3_4_15}
  \leanok
  \uses{def:pet-ch3-curvature-operator,def:pet-ch3-ricci-curvature}
  For a $3$-manifold, if the curvature operator in diagonal form is
  $\operatorname{diag}(\alpha,\beta,\gamma)$, show the Ricci curvature is diagonal
  with entries $\alpha+\beta,\beta+\gamma,\alpha+\gamma$, and that
  $\alpha,\beta,\gamma$ must be sectional curvatures.
  \source{petersen:page-0123}
\end{remark}

\begin{remark}[Exercise 3.4.16]
  \label{rem:pet-ch3-ex-16}
  \lean{PetersenLib.exercise3_4_16}
  \leanok
  \uses{def:pet-ch3-ricci-curvature,def:pet-ch3-scalar-curvature,notation:pet-ch3-divergence-adjoint}
  For the $(0,2)$-tensor $T=\operatorname{Ric}+b\cdot\operatorname{scal}\cdot
  g+cg$, $b,c\in\mathbb R$: (1) show $\nabla^*T=0$ if $b=-\tfrac12$ — the tensor
  $G=\operatorname{Ric}-\tfrac{\operatorname{scal}}2g+cg$ is the \emph{Einstein
  tensor}, $c$ the \emph{cosmological constant}; (2) if $c=0$, show $G=0$ in
  dimension $2$; (3) for $n>2$, show $G=0$ implies the metric is Einstein; (4)
  for $n>2$, show $G=0$ and $c=0$ implies the metric is Ricci flat.
  \source{petersen:page-0123,petersen:page-0124}
\end{remark}

\begin{remark}[Exercise 3.4.17]
  \label{rem:pet-ch3-ex-17}
  \lean{PetersenLib.exercise3_4_17}
  \leanok
  % Lean: parts (1),(2),(3) — the complexification's structural facts (isotropy
  % characterisation, isotropic-frame orthogonality with the √2-normalisation
  % g(v_i,v_i)=1/2 that Petersen calls "orthonormal", and reality of the complex
  % sectional curvature R_ℂ(v,w,w̄,v̄)). Parts (4),(5) (Berger quarter-pinching /
  % curvature-operator positivity) are not yet formalised.
  \uses{def:pet-ch3-sectional-curvature}
  Let $T_{\mathbb C}M=TM\otimes\mathbb C$; for $v=v_1+iv_2\in T_{\mathbb C}M$
  ($v_1,v_2\in TM$), $\bar v=v_1-iv_2$, and any tensorial object on $TM$
  complexifies, e.g.\ for a $(1,1)$-tensor $S$,
  $S_{\mathbb C}(v_1+iv_2)=S(v_1)+iS(v_2)$. A Riemannian $g$ induces a Hermitian
  structure $g(v,w)=g_{\mathbb C}(v,\bar w)=g(v_1,w_1)+g(v_2,w_2)+i(g(v_2,w_1)-g(v_1,w_2))$.
  A vector $v$ is \emph{isotropic} if $g(v,\bar v)=0$, i.e.\
  $g(v_1,v_1)=g(v_2,v_2)$ and $g(v_1,v_2)=0$; more generally isotropic subspaces
  are those on which $g_{\mathbb C}$ vanishes. The \emph{complex sectional
  curvature} of Hermitian-orthonormal $v,w$ is $R_{\mathbb C}(v,w,\bar w,\bar v)$,
  called \emph{isotropic} when $v,w$ span an isotropic plane. (1) Show $v_1+iv_2$
  is isotropic if $v_1,v_2$ are orthogonal of equal length. (2) Show that if an
  isotropic plane is spanned by Hermitian-orthonormal isotropic $v=v_1+iv_2$,
  $w=w_1+iw_2$, then $v_1,v_2,w_1,w_2$ are orthonormal. (3) Show
  $R_{\mathbb C}(v,w,\bar w,\bar v)$ is always real. (4) Show that if the metric
  is strictly quarter-pinched ($\sec\in(\tfrac14k,k)$, $k>0$), the complex
  sectional curvatures are positive. (5) Show the complex sectional curvatures
  are nonnegative (resp.\ positive) if the curvature operator is (resp.\
  strictly). Hint: compare
  $g(\mathfrak R(x\wedge u-y\wedge v),x\wedge u-y\wedge v)+g(\mathfrak R(x\wedge
  v+y\wedge u),x\wedge v+y\wedge u)$ to a suitable complex curvature.
  \source{petersen:page-0124,petersen:page-0125}
\end{remark}

\begin{remark}[Exercise 3.4.18]
  \label{rem:pet-ch3-ex-18}
  \lean{PetersenLib.exercise3_4_18}
  \leanok
  % Lean: parts (1),(3),(4) and the sec/scal halves of (2); the curvature-operator R halves of (2) omitted (follows from sec scaling).
  \uses{def:pet-ch3-sectional-curvature,def:pet-ch3-scalar-curvature,def:pet-ch3-curvature-operator,def:pet-ch3-ricci-curvature}
  Scaling $(M,g)$ to $(M,\lambda^2g)$: (1) show the connection and $(1,3)$-curvature
  tensor are unchanged; (2) $\sec,\operatorname{scal},\mathfrak R$ are all
  multiplied by $\lambda^{-2}$; (3) $\operatorname{Ric}$ as a $(1,1)$-tensor is
  multiplied by $\lambda^{-2}$; (4) $\operatorname{Ric}$ as a $(0,2)$-tensor is
  unchanged.
  \source{petersen:page-0125}
\end{remark}

\begin{remark}[Exercise 3.4.19]
  \label{rem:pet-ch3-ex-19}
  \lean{PetersenLib.exercise3_4_19}
  \leanok
  \uses{def:pet-ch3-curvature-tensor}
  Call $X$ an \emph{affine vector field} if $\mathcal L_X\nabla=0$. Show such a
  field satisfies $\nabla^2_{U,V}X=-R(X,U)V$.
  \source{petersen:page-0125}
\end{remark}

\begin{remark}[Exercise 3.4.20 — Integrability for PDEs]
  \label{rem:pet-ch3-ex-20}
  \lean{PetersenLib.exercise3_4_20}
  \leanok
  % Lean: part (1) only — the NECESSITY of the integrability conditions
  % (Ch03/PDEIntegrability.lean, exercise3_4_20).  Encoding P^i_k(x,u) as
  % P : E → F → (E →L F), the system reads fderiv u = fun y ↦ P y (u y); the
  % total x-derivative of the RHS is the chain-rule combination ∂₁P+(∂₂P)P, and
  % its symmetry in the two directions is exactly the integrability condition
  % ∂P^i_k/∂x^j+(∂P^i_k/∂u^l)P^l_j = ∂P^i_j/∂x^k+(∂P^i_j/∂u^l)P^l_k.  Proof:
  % because u solves the PDE, this object is the mixed second derivative D²u,
  % symmetric by Clairaut (second_derivative_symmetric_of_eventually) — no
  % smoothness of P is needed, a C² solution suffices.  The converse
  % sufficiency (part 2 = Frobenius theorem) and the flat-⟹-Cartesian-coords
  % consequences (parts 3–4) need integral-manifold / PDE-existence machinery
  % absent from Mathlib and PetersenLib, and are not formalized.
  \uses{prop:pet-ch3-curvature-coordinate-formula}
  For $P^i_k(x,u)$, $x=(x^1,\dots,x^n)$, $u=(u^1,\dots,u^m)$, consider
  $\partial u^i/\partial x^k=P^i_k(x,u(x))$, $u(x_0)=u_0$. (1) Show
  $\partial^2u^i/\partial x^k\partial x^j=\partial P^i_k/\partial
  x^j+(\partial P^i_k/\partial u^l)P^l_j$, and conclude solvability requires the
  \emph{integrability conditions}
  $\partial P^i_k/\partial x^j+(\partial P^i_k/\partial u^l)P^l_j=\partial
  P^i_j/\partial x^k+(\partial P^i_j/\partial u^l)P^l_k$. (2) Conversely show
  the integrability conditions suffice (equivalent to the Frobenius theorem). (3)
  For the system $\partial U^i_k/\partial x^j=\Gamma^i_{kj}U^i_j$ on a Riemannian
  $n$-manifold, show the integrability conditions are equivalent to
  $R^i_{kij}=0$. (4) Show a flat Riemannian manifold admits Cartesian
  coordinates, via $\partial u^i/\partial x^k=U^i_k$ with appropriate initial
  values, checking $g^{kl}(\partial u^i/\partial x^k)(\partial u^j/\partial x^l)$
  is constant using \cref{prop:pet-ch3-curvature-coordinate-formula}.
  \source{petersen:page-0125,petersen:page-0126}
\end{remark}

\begin{remark}[Exercise 3.4.21 — fundamental theorem of (hyper-)surface theory]
  \label{rem:pet-ch3-ex-21}
  \lean{PetersenLib.exercise3_4_21}
  \uses{prop:pet-ch3-curvature-coordinate-formula,thm:pet-ch3-tangential-curvature-equation,thm:pet-ch3-normal-curvature-equation}
  For a Riemannian immersion $F=(u^1,\dots,u^{n+1}):M^n\hookrightarrow\mathbb
  R^{n+1}$, $U^i_k=\partial u^i/\partial x^k$: (1) show
  $\partial U^i_j/\partial x^k=\Gamma^i_{kj}U^i_j-\Pi_{jk}N^i$ for a unit normal
  $N=N^i\partial/\partial u^i$ and $\Pi_{jk}=\Pi(\partial_j,\partial_k)=g(\nabla_jN,\partial_k)$;
  (2) show the integrability conditions are equivalent to the Gauss and Codazzi
  equations, $R_{ikij}=\Pi_{ij}\Pi_{kl}-\Pi_{ik}\Pi_{jl}$; (3) conversely, given
  $g_{ij}$ and symmetric $\Pi_{ij}$ satisfying Gauss--Codazzi, show a local
  Riemannian immersion with second fundamental form $\Pi_{ij}$ exists; (4) for a
  metric of constant curvature $R^{-2}>0$, show there is a local Riemannian
  immersion into $\mathbb R^{n+1}$ landing in a sphere of radius $R$ (with unit
  normal $N=\pm R^{-1}F$).
  \source{petersen:page-0126,petersen:page-0127}
\end{remark}

\begin{remark}[Exercise 3.4.22]
  \label{rem:pet-ch3-ex-22}
  \lean{PetersenLib.exercise3_4_22}
  \uses{rem:pet-ch3-ex-21}
  Repeat \cref{rem:pet-ch3-ex-21} for an immersion $F:M^n\hookrightarrow\mathbb
  R^{n,1}$ with normal $N$ satisfying $|N|^2=-1$, obtaining a local
  characterization of the hyperbolic spaces $H^n(R)$ as the local model of
  constant curvature $-R^{-2}$ (with unit normal $N=\pm R^{-1}F$).
  \source{petersen:page-0127}
\end{remark}

\begin{remark}[Exercise 3.4.23 — Kulkarni--Nomizu product]
  \label{rem:pet-ch3-ex-23}
  \lean{PetersenLib.exercise3_4_23}
  \leanok
  \uses{prop:pet-ch3-curvature-symmetries,def:pet-ch3-constant-curvature}
  For symmetric $(0,2)$-tensors $h,k$, define
  $h\owedge k(v_1,v_2,v_3,v_4)=\tfrac12(h(v_1,v_4)k(v_2,v_3)+h(v_2,v_3)k(v_1,v_4))
  -\tfrac12(h(v_1,v_3)k(v_2,v_4)+h(v_2,v_4)k(v_1,v_3))$. (1) Show $h\owedge
  k=k\owedge h$. (2) Show $h\owedge h=0$ if $h$ has rank $1$. (3) Show that if
  $n>2$, $k$ is nondegenerate, and $h\owedge k=0$, then $h=0$. (4) Show $h\owedge
  k$ satisfies properties (1)--(3) of \cref{prop:pet-ch3-curvature-symmetries}.
  (5) Show $\nabla_X(h\owedge k)=(\nabla_Xh)\owedge k+h\owedge(\nabla_Xk)$. (6)
  Show $(M,g)$ has constant curvature $c$ iff the $(0,4)$-curvature tensor
  satisfies $R=c\cdot(g\owedge g)$.
  \source{petersen:page-0127}
\end{remark}

\begin{remark}[Exercise 3.4.24 — Schouten tensor]
  \label{rem:pet-ch3-ex-24}
  \lean{PetersenLib.exercise3_4_24}
  \leanok
  \uses{def:pet-ch3-ricci-curvature,def:pet-ch3-scalar-curvature,rem:pet-ch3-ex-23}
  Define $P=\dfrac2{n-2}\operatorname{Ric}-\dfrac{\operatorname{scal}}{(n-1)(n-2)}g$
  for $n>2$. (1) Show $P\equiv0$ implies $\operatorname{Ric}=0$. (2) Show the
  decomposition
  $P=\dfrac{\operatorname{scal}}{n(n-1)}g+\dfrac2{n-2}(\operatorname{Ric}-\tfrac{\operatorname{scal}}ng)$
  is orthogonal. (3) For $n=2$, show $R=\tfrac{\operatorname{scal}}2g\owedge g$.
  (4) For $n=3$, show
  $R=\tfrac{\operatorname{scal}}6g\owedge g+2(\operatorname{Ric}-\tfrac{\operatorname{scal}}3g)\owedge
  g=P\owedge g$. (5) For $n>2$, show $(M,g)$ has constant curvature iff
  $R=P\owedge g$ and $\operatorname{Ric}=\tfrac{\operatorname{scal}}ng$. (6) Show
  $\operatorname{Ric}(X,Y)=\sum_i(P\owedge g)(X,E_i,E_i,Y)$ for any orthonormal
  frame $E_i$.
  \source{petersen:page-0127,petersen:page-0128}
\end{remark}

\begin{remark}[Exercise 3.4.25 — Weyl tensor]
  \label{rem:pet-ch3-ex-25}
  \lean{PetersenLib.exercise3_4_25}
  \leanok
  \uses{rem:pet-ch3-ex-24}
  Define $W$ implicitly by
  $R=\dfrac{\operatorname{scal}}{n(n-1)}g\owedge g+\dfrac2{n-2}(\operatorname{Ric}-\tfrac{\operatorname{scal}}ng)\owedge
  g+W=P\owedge g+W$ (with $P$ as in \cref{rem:pet-ch3-ex-24}). (1) Show $n=3$
  implies $W=0$. (2) Show $\sum_iW(X,E_i,E_i,Y)=0$ for any orthonormal frame
  $E_i$. (3) Show the decomposition $R=P\owedge g+W$ is orthogonal.
  \source{petersen:page-0128}
\end{remark}

\begin{remark}[Exercise 3.4.26]
  \label{rem:pet-ch3-ex-26}
  \lean{PetersenLib.exercise3_4_26}
  \leanok
  % Lean: first identity only (div-adjoint of the Schouten tensor = -1/(n-1) d(scal)); the Weyl-divergence identity needs (0,4)-tensor divergence infra (not yet formalized).
  \uses{rem:pet-ch3-ex-24,notation:pet-ch3-divergence-adjoint}
  Show $\nabla^*P=-\dfrac1{n-1}d\operatorname{scal}$ and
  $\nabla^*W(Z,X,Y)=\dfrac{n-3}2\big((\nabla_XP)(Y,Z)-(\nabla_YP)(X,Z)\big)$.
  \source{petersen:page-0129}
\end{remark}

\begin{remark}[Exercise 3.4.27 — structure constants]
  \label{rem:pet-ch3-ex-27}
  \lean{PetersenLib.exercise3_4_27}
  \leanok
  \uses{def:pet-ch3-curvature-tensor}
  For an orthonormal frame $E_1,\dots,E_n$ define structure constants $c^k_{ij}$
  by $[E_i,E_j]=c^k_{ij}E_k$ (functions on $M$), and $\Gamma^k_{ij},R^l_{ijk}$ by
  $\nabla_{E_i}E_j=\Gamma^k_{ij}E_k$, $R(E_i,E_j)E_k=R^l_{ijk}E_l$. Compute
  $\Gamma,R$ in terms of the $c^k_{ij}$; note that on a Lie group with a
  left-invariant metric the $c^k_{ij}$ may be taken constant.
  \source{petersen:page-0129}
\end{remark}

\begin{remark}[Exercise 3.4.28 — Cartan formalism]
  \label{rem:pet-ch3-ex-28}
  \lean{PetersenLib.exercise3_4_28}
  \leanok
  % Lean: parts (1)-(2) formalized for a smooth orthonormal frame -- skew-symmetry
  % $\omega^i_j=-\omega^j_i$ (connectionForm_skew), the first structure equation
  % $d\omega^i=\omega^j\wedge\omega^i_j$ (firstStructureEquation), and the second
  % $d\omega^j_i=\omega^k_i\wedge\omega^j_k+\Omega^j_i$ (secondStructureEquation),
  % using the intrinsic $d\theta(X,Y)=X\theta(Y)-Y\theta(X)-\theta([X,Y])$ and a
  % frame Parseval identity (sum_metricInner_frame). The converse (that (1)
  % determines the connection forms) and the surface case (3), whose
  % $\Omega^1_2=\sec\cdot d\mathrm{vol}$ needs a Riemannian area form, are not formalized.
  \uses{def:pet-ch3-curvature-tensor}
  For a frame $E_1,\dots,E_n$, write the connection as $\nabla E_i=\omega^j_iE_j$
  (connection $1$-forms $\omega^j_i$), i.e.\ $\nabla_vE_i=\omega^j_i(v)E_j$; for
  an orthonormal frame with dual coframe $\omega^i$: (1) show $\omega^i_j=-\omega^j_i$
  and $d\omega^i=\omega^j\wedge\omega^i_j$, and that these two equations
  conversely determine the connection forms from the orthonormal frame. (2)
  Viewing $\omega^j_i$ as an $\mathfrak{so}(n)$-valued $1$-form, define the
  curvature forms $\Omega^j_i$ ($\mathfrak{so}(n)$-valued $2$-forms) by
  $R(X,Y)E_i=\Omega^j_i(X,Y)E_j$; show
  $d\omega^j_i=\omega^k_i\wedge\omega^j_k+\Omega^j_i$. (3) For surfaces, with
  orthonormal frame $E_1,E_2$ and coframe $\omega^1,\omega^2$, show
  $d\omega^1=\omega^2\wedge\omega^1_2$, $d\omega^2=-\omega^1\wedge\omega^1_2$,
  $d\omega^1_2=\Omega^1_2=\sec\cdot d\mathrm{vol}$.
  \source{petersen:page-0129,petersen:page-0130}
\end{remark}

\begin{remark}[Exercise 3.4.29]
  \label{rem:pet-ch3-ex-29}
  \lean{PetersenLib.exercise3_4_29}
  \leanok
  \uses{lem:pet-ch3-bivector-jacobi-identity,prop:pet-ch3-curvature-symmetries}
  Assume $R(X,Y,Z,W)$ is an algebraic curvature tensor, i.e.\ satisfies
  properties (1)--(3) of \cref{prop:pet-ch3-curvature-symmetries}. (1) Show
  \begin{align*}
    6R(X,Y,V,W)&=\frac{\partial^2}{\partial s\,\partial t}R(X+sW,Y+tV,Y+tW,X+sW)\Big|_{s=t=0}\\
    &\quad-\frac{\partial^2}{\partial s\,\partial t}R(X+sV,Y+tW,Y+tW,X+sV)\Big|_{s=t=0}.
  \end{align*}
  (2) Show $6R(X,Y,V,W)$ equals the explicit sum-of-twelve-terms expansion
  \begin{align*}
    &R(X+W,Y+V,Y+V,X+W)-R(X,Y+V,Y+V,X)-R(W,Y+V,Y+V,W)\\
    &\quad-R(X+W,V,V,X+W)-R(X+W,Y,Y,X+W)+R(X,V,V,X)+R(W,V,V,W)\\
    &\quad+R(X,Y,Y,X)+R(W,Y,Y,W)-R(X+V,Y+W,Y+W,X+V)\\
    &\quad+R(X,Y+W,Y+W,X)+R(V,Y+W,Y+W,V)+R(X+V,Y,Y,X+V)+R(X+V,W,W,X+V)\\
    &\quad-R(X,Y,Y,X)-R(V,Y,Y,V)-R(X,W,W,X)-R(V,W,W,V)
  \end{align*}
  (four terms are redundant, obtainable from polarizing $R(X,Y,Y,X)$).
  \source{petersen:page-0130,petersen:page-0131}
\end{remark}

\begin{remark}[Exercise 3.4.30]
  \label{rem:pet-ch3-ex-30}
  \lean{PetersenLib.exercise3_4_30}
  \leanok
  \uses{rem:pet-ch3-ex-29,def:pet-ch3-curvature-operator}
  Using \cref{rem:pet-ch3-ex-29}, show the norm of the curvature operator on
  $\Lambda^2T_pM$ is bounded, $|\mathfrak R|_p\le c(n)|\sec|_p$, for a constant
  $c(n)$ depending only on dimension, where $|\sec|_p$ is the largest absolute
  value of a sectional curvature at $p$.
  \source{petersen:page-0131}
\end{remark}

\begin{remark}[Exercise 3.4.31]
  \label{rem:pet-ch3-ex-31}
  \lean{PetersenLib.exercise3_4_31}
  \leanok
  \uses{def:pet-ch3-curvature-tensor}
  Let $G$ be a Lie group with a left-invariant nondegenerate metric
  $(\cdot,\cdot)$ on $\mathfrak g$; for $X\in\mathfrak g$ let
  $\operatorname{ad}_X^*$ be the adjoint of $\operatorname{ad}_XY=[X,Y]$ with
  respect to $(\cdot,\cdot)$, so $(\operatorname{ad}_X^*Y,Z)=(Y,[X,Z])$. For
  left-invariant fields:
  (1) the Levi-Civita connection is
  $\nabla_XY=\tfrac12([X,Y]-\operatorname{ad}_X^*Y-\operatorname{ad}_Y^*X)\in\mathfrak g$;
  it is torsion-free and metric-compatible.
  (2) $R(X,Y,Z,W)=-(\nabla_YZ,\nabla_XW)+(\nabla_XZ,\nabla_YW)-(\nabla_{[X,Y]}Z,W)$.
  (3)
  $R(X,Y,Y,X)=\tfrac14|\operatorname{ad}_X^*Y+\operatorname{ad}_Y^*X|^2-(\operatorname{ad}_X^*X,\operatorname{ad}_Y^*Y)-\tfrac34|[X,Y]|^2-\tfrac12([[X,Y],Y],X)-\tfrac12([[Y,X],X],Y)$.
  We treat the positive-definite (Riemannian) case. (Petersen's printed part~(1)
  has $+\operatorname{ad}_X^*Y$; with the adjoint convention above the correct
  sign is $-\operatorname{ad}_X^*Y$, as the biinvariant specialisation
  $\nabla_XY=\tfrac12[X,Y]$ confirms.)
  \source{petersen:page-0131}
\end{remark}

\begin{remark}[Exercise 3.4.32]
  \label{rem:pet-ch3-ex-32}
  \lean{PetersenLib.exercise3_4_32}
  \leanok
  \uses{rem:pet-ch3-ex-31,def:pet-ch3-curvature-operator}
  Let $G$ be a Lie group with a biinvariant (possibly indefinite, nondegenerate)
  metric $(\cdot,\cdot)$ on $\mathfrak g$. Using left-invariant fields, show: (1)
  $\nabla_XY=\tfrac12[X,Y]$; (2) $R(X,Y)Z=\tfrac14[Z,[X,Y]]$; (3)
  $R(X,Y,Z,W)=-\tfrac14([X,Y],[Z,W])$, so sectional curvatures are nonnegative
  when $(\cdot,\cdot)$ is positive definite; (4) the curvature operator is
  nonnegative when $(\cdot,\cdot)$ is positive definite, by showing
  $g\big(\mathfrak R(\sum_{i=1}^kX_i\wedge Y_i),\sum_{i=1}^kX_i\wedge
  Y_i\big)=\tfrac14\sum_{i=1}^k|[X_i,Y_i]|^2$; (5) for $(\cdot,\cdot)$ positive
  definite, $\operatorname{Ric}(X,X)=0$ iff $X$ commutes with every left-invariant
  vector field, so $G$ has positive Ricci curvature if its center is discrete.
  \source{petersen:page-0131,petersen:page-0132}
\end{remark}

\begin{remark}[Exercise 3.4.33]
  \label{rem:pet-ch3-ex-33}
  \lean{PetersenLib.exercise3_4_33}
  \leanok
  \uses{rem:pet-ch3-ex-32,def:pet-ch3-ricci-curvature}
  For a Lie group with nondegenerate Killing form $B$, using $-B$ as the
  left-invariant metric: (1) show this metric is biinvariant; (2) show
  $\operatorname{Ric}=-\tfrac14B$.
  \source{petersen:page-0132}
\end{remark}

\begin{remark}[Exercise 3.4.34]
  \label{rem:pet-ch3-ex-34}
  \lean{PetersenLib.exercise3_4_34}
  \leanok
  \uses{rem:pet-ch3-ex-33}
  Using the Cartan formalism of \cref{rem:pet-ch3-ex-28} in the setting of
  \cref{rem:pet-ch3-ex-33}, compute all quantities in terms of the structure
  constants of $\mathfrak g$: with a biinvariant metric, in an orthonormal basis
  they satisfy $c^k_{ij}=-c^k_{ji}=c^i_{jk}$, the first equality being
  skew-symmetry of the bracket, the second biinvariance of the metric.
  \source{petersen:page-0132}
\end{remark}
